\documentclass[11pt,a4paper]{article}
\usepackage[margin=22mm]{geometry}
\usepackage{amsmath,amssymb,amsthm,booktabs,longtable,array,fontspec}
\setmainfont{Noto Sans CJK SC}[AutoFakeBold=2,AutoFakeSlant=0.15]
\setmonofont{lmmono10-regular.otf}[Scale=0.85]
\XeTeXlinebreaklocale "zh"
\XeTeXlinebreakskip=0pt plus 1pt
\usepackage[colorlinks=true,linkcolor=blue,urlcolor=blue]{hyperref}
\usepackage{xurl,fvextra}
\DefineVerbatimEnvironment{verbatim}{Verbatim}{fontsize=\small,breaklines=true}
\newcommand{\tightlist}{\setlength{\itemsep}{0pt}\setlength{\parskip}{0pt}}
\newcommand{\FS}{\operatorname{FS}}
\newcommand{\Z}{\mathbb Z}
\newcommand{\A}{A_*}
\newcommand{\eps}{\varepsilon}
\newtheorem{theorem}{定理}[section]
\newtheorem{lemma}[theorem]{引理}
\renewcommand{\proofname}{证明}
\renewcommand{\abstractname}{摘要}
\setlength{\emergencystretch}{3em}
\renewcommand{\arraystretch}{1.12}
\linespread{1.04}\raggedbottom
\title{排除 84 元配置，将上界从 84 改进到 83\\\large 完整结构归约、实际资源证明与精确有限核验}
\author{Sum-of-Subset 研究项目}
\date{2026 年 10 月 19 日}
\begin{document}
\maketitle
\begin{abstract}
从提交 \texttt{cf26c50} 已独立闭合的上界 84 出发，排除大小恰为 84 的可行集合，得到
$71\le N\le83$。本次不再只给出条件标量系统：重新计算大小 84 的共同整数源归约，
处理等差支撑的边界，给出全部高洞情形的区间延拓；低正洞统一局部化，随后按空覆盖单元数处理。
主要新步骤是两空单元的统一证明、四空无所有者族的两次截断，以及五空 alpha 族
先证明次序、后证明正常性的扣减链 $14\to27\to31\to35\to38\to42$。
五空正洞只需微小剩余预算。报告明确区分无限集合论证、有限参数验证和组件分类。
验证器使用 Python 标准库及整数或有理数运算；它不是形式化证明助手。
\end{abstract}
\tableofcontents
\newpage

\section{结论、前提与证明接口}
\begin{theorem}\label{thm:main}
若有限实数集 $A$ 满足 $\FS(A)\subseteq A+A+A+A$，则 $|A|\le83$。
结合仓库已验证的 71 元构造，有 $\boxed{71\le N\le83}$。
\end{theorem}
子集和包括空集和，四项和允许重复。已有上界 84 作为归纳前提；若 84 元反例不含零，
加入零会产生 85 元反例。因此可假设零属于反例。
为每个子集和选定一个四项表示，这些表示构成有限个齐次整数线性方程。
其有理核在实核中稠密，连同指定坐标为零的方程，可取保持所有严格次序及符号的有理解，
再乘正公分母。因此只需排除一个含零的 84 元整数反例；并未断言每个原实数集本身有公度。

本报告沿用已经逐项核查的通用引理：四极值覆盖与剥离、非等差缺陷递推、计数洞引理、
有效洞引理、同时逆定理及短直径定理、整数四极值限制和集界、实际资源分组损失引理。
依赖源文本包括旧完整报告的附录 D/E、\texttt{STRUCTURAL86.md} 所列通用定理，
以及大小 85 的结构伴随证明；本次仅复用其通用论证，所有与基数有关的有限参数重新计算。
完整路径及 SHA-256 见 \path{data/upper83_complete_sources.json}。
旧上界 84 的 140 项有限义务由原完整入口重新运行；不能以旧 85 元分类表代替新 84 元分类。

\subsection{记号与实际覆盖}
去掉两侧各四个外点，得 $|D|=76$；再做两次两侧各两点的剥离，得
\[
 E=P_s\cup\{0\}\cup(-P_t),\quad |E|=68,\quad r_s+r_t=67.
\]
标量分析可反射为 $r_s=33-k,r_t=34+k$，但所有者证明始终保留物理符号。
令 $T_r=r(r+1)/2$，$I=k(k+1)=(r_v-33)(r_v-34)$。
核心点写成 $x_i=i+d_i$，$d_i\ge0$ 不减；定义
\[
 \tau_v=\sum d_i,\quad \xi_v=d_{r_v},\quad R_v=r_v+\xi_v,
 \quad R=67+\xi,\quad 0\le\xi\le\tau.
\]
每侧四个源尾点为 $R_v+a_v,R_v+b_v,R_v+c_v,R_v+e_v$，偏移为严格递增正整数。
令 $M_v=R_v+e_v$，$L=e_s+e_t$，$M=R+L$，$U=\sum_v(a_v+b_v)$，
$C=c_s+c_t$，$J=L-C$。$S_v$ 是 $D_v$ 正点总和，其中 $D_s=D,D_t=-D$。
外点为 $v_1>v_2>v_3>v_4>M_v$，定义
\[
 d_v=v_1-v_2,\ g_v=v_3-v_4,\ h_v=v_2-v_3,\ X_v=v_4-M_v,
 \quad t_v=d_v-M,
\]
\[
 \lambda_v=d_v-v_3-v_4,\quad
 U_v=d_v-v_4,\ V_v=d_v-v_3,\ H_v=d_v+h_v-v_4,\ O_v=H_v+d_v,
\]
\[
 \alpha_v=h_v-v_4,\quad\beta_v=\alpha_v-d_v,
 \quad C_{iv}=\lambda_v-o_{i-3}\ (4\le i\le7).
\]
这里普通中心 $U_v$ 与源偏移和 $U$ 不同，固定标签 $U_i^v$ 又是另一对象。

两个首项资源为 $\lambda_v+(D_v+D_v)$，容量 $K=|D+D|$。
普通覆盖有 18 个整单元，每个容量 76：$C_4$ 始终配对；
分支 A 在 $\beta_s+\beta_t\le0$ 时配对 beta，分支 B 配对 $C_5$。
一个配对只拥有一次容量，只有实际唯一贡献侧才能按原生普通资源计损失。
固定 110 个标签中的 26 个互不相交负和组给出固定容量 84。
若有 $j$ 个空整单元，令 $c=76j+\delta$，则
\begin{equation}\label{eq:source-mass}
 S_s+S_t=2K+(18-j)76+84+w-\delta,\qquad \delta\ge0.
\end{equation}
$w$ 是 $\FS(D)$ 的整数凸包洞数；两物理目标的洞数之和是 $w$。
对互不相交的实际资源组，其容量减去目标并集的损失之和不超过 $\delta$。
相容额外固定扣减 $e$ 将预算降为 $\eps=\delta-e$。

\section{共同格子与结构归约}
定义
\[
 \mu_X=|\FS(X)|-1-T_{|X|},\quad
 \rho_X=\mu_X+2-2\Delta_0(X),
\]
其中 $\Delta_0(X)$ 数 $X$ 中不能在 $X\cup\{0\}$ 中表示为不同元素和的对角点。
令 $\kappa=|\FS(P_s)+\FS(P_t)|-|\FS(P_s)|-|\FS(P_t)|+1$。
通用剥离恒等式在本基数成为
\[
 \Lambda=1157+I,\quad\Omega=612-I,
 \quad\rho+\kappa+2\theta+c+\ell_1+\ell_2+g=\Omega.
\]
这里所有松弛非负，$\rho$ 在小边界允许已知下界 $-8$。
两侧符号不能消失：初始单侧下界 $1+T_{68}=2347>2093$，
后续下界 $1+T_{67}=2279>1901$。未配对点满足 $q^2\le2977$，奇数 $q\le53$，
故至少七对相反点；$\Omega\ge-8$ 进一步给较小侧至少九点。

获得共同核心格子以后、尚未使用整数四极值不等式之前，
\[
 |\FS(D_1)|-2|\FS(E)|\le-125-2I+2z,
\]
小侧不超过 14 时右侧加 16。负上界同时强制下一层同格和两极端平移相交。
外层由 $\Lambda>294+\Omega-\rho$ 强制同样结论。
完成这两步之后才能使用 $G=2\cdot76-10=142$，得到
\begin{equation}\label{eq:structural}
 \mathcal T\ge\max(48,z,142-2z),\quad
 \mathcal T\ge\kappa+142-3z,\quad c+\rho+\mathcal T\le612-I,
 \quad w\le4z.
\end{equation}
$h_s,h_t\le\kappa$，$z=h_s+h_t-\kappa\le\min(h_s,h_t)$。
以下是完整归约的关键修复，详细参数和边界在附录 A。

\begin{enumerate}
\item 较大支撑为 AP、较小支撑非 AP 时，使用 $\rho_s\ge n(r_s)$ 后得到
$\kappa<T_{r_t}$，强制单一剩余类和 $z=0$。两支撑都为 AP 时，
较小侧 $a\le32$ 仍有 $\Omega-\rho_s<T_b$；仅 $a=33,b=34$ 需另论。
此时互素内禀步长只有 $(1,1),(1,2),(2,1)$，不等步长导致 $z=0$、
$\kappa\ge T_{33}$，由后述小 $c$ 引理排除。
不能在所有不平衡两 AP 行直接声称 $\Omega<2T_a$。
\item 仅较小支撑 AP 时，正确阈值是
$\kappa\le\Omega-\rho_s-(b-5)<2T_a$。
单一剩余类的截断长洞或唯一离格生成元代价超过预算。
两剩余类是二阶剩余类；规范为偶数 AP 后，最小奇生成元 $W$ 的质量下界
\[
 \mu_t+\kappa\ge T_a+1+(b-t)(b-t+1)+tW+t(t-1)-W-T_b
\]
在 $W=19$ 已超过允许量，且对 $W$ 斜率 $t-1>0$。
故 $W\le17,z\le16$，格化后 $c\le93<152,w\le64$。
非本原单侧 AP 的相应界为 $c\le85<152$。
\item 两支撑非 AP 时，较大侧严格逆条件自动成立；较小侧充分条件仅 $k=0,1$
自动成立。实际不成立的 $k=2,\ldots,21$ 使用新计算的 20 个计数端点修复。
每行核对起点、离格、后继区间拼接和短直径的严格余量；主范围修复后 $z\le18$。
较小非本原侧给 $c\le65<76,w\le72$，由全非空引理排除。
\item 实际严格逆条件成立时，$z\le62+I$ 已使内层比较严格；其余 $z\ge63+I$
使用与基数无关的 $h\ge49$ 轮廓。若一侧最后两平移不交，其代价
$T_{r-2}+2-2r$ 加另一侧轮廓与 $63+I$ 超过 $\Omega$，38 个定向行最小余量为 77。
于是两侧末次平移都相交，最大元删除将内层比较加强为 $-125-2I<0$。
这里未提前使用 $G=142$。
\end{enumerate}
$a=9,\ldots,14$ 的 AP、实际严格逆和端点修复边界全部由附录 A 分别排除。
最终获得两侧本原的共同整数源，$15\le r_s,r_t\le52$。
若 $z>0$，两侧支撑均非 AP，故
\begin{equation}\label{eq:positivefloor}
 \mu\ge61,\quad \tau=\mu+h_s+h_t\ge61+2z.
\end{equation}

\subsection{不需要本原性的低容量例外引理}
通用守护空隙计数在全非空时给 $w\ge2M+10\ge160$。
对上述 $c<152,w\le64$ 的结构例外，若 $j=0$ 即矛盾；若 $j=1$，$\delta\le75$。
一个空整单元的原生中心恒等式确定整性，实际 $C_6$ 或 $C_7$ 给 $d_v\ge80$。
原始帽形界 $t_v\le12+\eta_v\le76<80$、$2t_v\le17+\eta_v\le81$
足以得到唯一所有者与六边森林。源界给 $\tau\le444$；
全秩核中若 $\phi_s+\phi_t\ge42$，质量至少 $2T_{21}=462$，故 $F\ge27$。
于是 $2X+G\ge6F-\delta\ge87$，但所有者恒等式给
$2X+G\le(17+64)/2-2=38.5$。此处只用 $\eta_v\le w$，未提前假定对称洞数。

\section{源包络、实际截断与区间延拓}
令 $\A=I+\tau+\delta$。通用源定理给
\[
 \Theta=76j-296-w+\A\le160,\qquad\Theta+B_s+B_t\le170,
\]
其中 $B_v$ 是四个源偏移两两和超过最大偏移的实际不同值数。
结合删除一或两个最大源点的包络，得到
\begin{align}
 \A&\le456-76j+w,& L&\ge\Theta-8,\\
 J+U&\le328-\Theta,&J&\ge\Theta-18,\\
 K&\ge\Theta+2R+2,&K-M&\ge R+76j-454-w+\A,\\
 \max(4,\Theta-88)&\le e_v\le325-\Theta.&&
\end{align}
$j=5,w=0,\delta\ge10$ 时另有尖锐加强 $L\ge\A+78$。
这些源结论在先前全秩报告中已有完整证明；本次用相同实际定义。

对每一侧选择某个真实前缀，令
\[
 \phi_v=\min_j(j+d_{r_v-j}),\qquad F=68-\phi_s-\phi_t.
\]
所得同一个实际子核心 $Q$ 在整数位移 $x$ 下满足
$|Q\cap(Q+x)|\ge F-x$。按所需的 $\beta F+q_0|Q|$ 选择一个优化前缀，
而不能独立优化 $F$ 与 $|Q|$。逆阶梯质量公式给
\[
 \tau+F\ge68,\ \tau+2F\ge134,\ \tau+3F\ge198,
 \quad\tau+4F\ge260,\quad\tau+6F\ge378,
\]
\[
 2\tau-2\xi+7F\ge444,\qquad2\tau+6F\ge396.
\]
森林按实际目标并集计损失；若一个分离臂中有空顶点，其容量为零。
所有未使用的原生损失返还预算，绝不向空配对指定所有者。

\subsection{所有高洞 $z\ge15$}
令 $H=h_s+h_t=\kappa+z=\tau-\mu$，核心中央区间长度参数
$W_0=1156+I+\mu-H$。只要两侧核心质量都大于 $H$，
简单和集洞引理给中央区间 $[H,C_0-H]$，$C_0=1156+I+\tau$。
新的外层恒等式和两最大元源包络给
\[
 M\le836+H+(\rho-\mu)/2+z-I-c,
\]
\begin{equation}
 W_0-M\ge-904+4I+3c+\tfrac32(\mu+\rho)-3z.
\end{equation}
完全重放的通用稠密轮廓提供以下成对下界；$H$ 上界还应减 $I+c$。
\begin{center}
\begin{tabular}{rrrrrr}\toprule
$z$&$\mu$&$\rho$&$I+c$ 上界&$H$ 上界&$W_0-M$ 下界\\\midrule
15--16&418&288&214&246&107\\
17--20&461&331&179&219&224\\
21--48&503&373&191&289&266\\\bottomrule
\end{tabular}
\end{center}
由 $\min_r(T_r+I)=385$，三段各侧核心质量超过 $H$ 的余量至少为139、166、96。
在 $z\ge49$，$\mu\ge539,\rho\ge409$，$\kappa\le203-I-c$、
$H\le406-2I-2c$，各侧质量余量至少56。
若 $j\ge1$，$c\ge76$，宽度余量至少 $365+7I+6(c-76)>0$。
若 $j=0$，原始 $w\ge2M+10$ 与 $w\le4z$ 给 $M\le401$，而 $W_0\ge1289$。
因此八个源尾点逐个加入时，中央区间均能拼接；共同边界洞定理给偶数 $w\le z$。
$j\ge1$ 时 $\A\le380+w$ 与 $\tau\ge418+2z$ 矛盾；
$j=0,z\le48$ 与 $w\ge160$ 矛盾，尾段则由 $\tau\ge539+2z>456+w$ 排除。

\subsection{低正洞的统一局部化}
对 $1\le z\le14,j\ge1$，先只用 $w\le4z\le56$。
$\min(T_r+I)=385$、$\max(r-I)=34$ 和 $\mu\ge61$ 给
\[
 C_v-H\ge76j-w-10+\delta>0,
\]
\[
 W_0\ge822+76j-w+2I+\delta\ge842,
 \quad M_v\le655-76j+w-I-\delta\le635.
\]
故同样逐个延拓八个源尾点，得到偶数 $w\le z\le14$、$\eta_s=\eta_t=w/2$，
所有目标洞在末端边界条带。$j=0$ 和 $j\ge6$ 已分别由全非空界和正质量下界排除。

\section{从活动判据到完整终端}
零洞用 $Z_j$，低正洞用 $P_j$ 表示。
守护空隙的三个实际帽形界为
\[
 O_v\text{ 活动}:t_v\le12+\eta_v,\quad
 \beta_v,H_v\text{ 活动}:t_v\le9+\eta_v,\quad
 \beta_v,O_v\text{ 活动}:2t_v\le17+\eta_v.
\]
它们只在已证明的正首项最大值或相应普通中心守护条件下使用。
完整目标的整数活动区间为 $[1-M_v,S_v]$；正洞的相同区间由边界条带、
至少 18 个相反核心像与至多七个目标洞证明。
附录 B--H 逐一给出整性、活动链、资源名单与终端，以下突出新步骤。

\subsection{一空、两空及六空}
$Z_1,P_1$ 的实际所有者语法留下 beta、$H$、$O$ 同侧活动。
带首项根损失的六边森林给
\[
 2t_u\ge I+\max(99-5\xi+\tau,27+2\xi)-w\ge51+I-w,
\]
与 $2t_u\le17+w/2$ 矛盾，低正洞余量至少13。

$Z_2$ 先补足可能空首项的情形：若一首项空，$\delta\ge K$ 与
$K-\delta\ge I+\tau-8$ 强制 $I+\tau\le8$。
剩余普通损失不超过8，整性恒等式使它们正常。
空首项若 HIGH 至少两个普通资源各损失 $r_v+5>8$；若 LOW，
跨侧链强制分支 B 恰好空 $C_{6v},C_{7v}$，正常 $V_v$ 给 $X_v\ge68$，
相反活动 $C_7$ 与 $O$ 的帽形界矛盾。因此两首项都活动。

$P_2$ 另有 $\tau-\xi\ge32$：第一个绝对核心子集和洞小于最大核心点 $\le352$。
至少34点的一侧删去最大点后质量至少 $T_{33}=561$，其前缀仍缺这个洞。
任意 $n$ 个不同正整数若子集和不是完整区间，其质量超出 $T_n$ 至少 $n-1$。
用首次贪心失败及 $(n-k)(T_{k-1}+1)\ge n-1$ 即证。

假定 beta、$O$ 同侧活动，另一个 $O$ 必空。按第二空单元是否为 $V_v,C_{7v},C_{7o}$
分四类，实际有根森林及未用损失分别给
\[
 2t_v\ge43-w+I+\xi;\quad
 t_v\ge14-w+(\tau-\xi);\quad t_v\ge76-w;\quad t_v\ge126-w.
\]
第一式要求 $26\le3w/2\le21$；第二式在零洞至少14，在正洞至少32；
后两式也均超过帽形上界。剩余只有 A 中
$\{O_v,B\},\{O_s,O_t\},\{O_v,H_v\}$。
全部 alpha、$C_5,U,V,C_7$ 保留，实际所有者给 $2X+G\le10+w/2$，
六边森林给 $2X+G\ge6F-\delta$。因此
\[
 \A\ge368+I-w/2>304+w\ge\A,
 \qquad64+I\le3w/2\le21
\]
矛盾。这统一闭合 $Z_2,P_2$。

$Z_6$ 的源等号强制 $I=\tau=\delta=0$、$\Theta=160$。
两个首项满容量，$\lambda_v\ge2M_o+1$。
由 $J\ge142,U+J\le168,C\ge6$，有 $8M+4-S\ge120$。
两侧固定 $C_{14}$ 之和超过 $8M+4>S$，至少一侧四个固定顶标签全 HIGH，
产生正扣减，与 $\delta=0$ 矛盾。$j>6$ 的零洞预算直接为负。

\subsection{三空与四空的实际接入}
四顶尾部与四个完整核心星形给“两帽”必要条件
\[
 4I+4\tau+\delta+228j\le538+6w.
\]
它排除三空的两帽。四空在整性之前已可由 $2Y>482$ 与 $Y\le190+3w\le202$ 排除两帽。
随后两个见证的剩余类恒等式、四边下界 $248-\tau-\delta$、
五边下界 $300-\tau-\delta$ 强制外点整性。
特别是 HIGH 跨中心的计数要保留十个正正固定 HIGH 标签：三空 $C_5\Rightarrow C_6$
要求 $\A\ge I+4r+105\ge235$；四空末链要求 $\A\ge I+2r+95\ge161$。
这两行分别超过正洞预算234和156。

$j\ge3$ 的正洞核心含 $[1,17]$ 中至少12个不同值：六个同时缺值至少花费270质量。
6188个十二值种子的精确子集和检验给中央区间；其他核心点全部能顺序拼接。
故低洞只在 $[1,10]$。每值重数为0、1、2的59049种模式完整枚举给：
四个低洞成本至少256，三个至少195，第一个洞为2至少99。
由此 $P_3:w\le6,\A\le234$，$P_4:w\le4,\A\le156$，
$P_5:w=z=2,\tau=67+\sigma,I+\sigma+\delta\le11$。

三空实际分类只有 A0、A1；A0 的五个实际端点产生十二个重截断模式，
零洞最小余量65，正洞41。A1 的有根并集式给统一下界 $391.5$，
四个端点模式的零洞上界至多356、正洞至多374。
四空实际分类为 A-E、A-H、B。七个实际端点母式的最小余量为
\begin{center}
\begin{tabular}{lrrrrrrr}\toprule
母式&F&S&E-H&E-alpha自&E-alpha对&E-C5&O\\\midrule
零洞&306&81&253&163&204&295&123\\
正洞&290&65&237&151&164&283&103\\\bottomrule
\end{tabular}
\end{center}
所有由 $e+72h\le\A$ 允许的重截断标志均纳入，包括两个重端点。
O 母式的常数994从原生端点恒等式和三条森林相加推导，未从旧基数外推。

仅余四空无所有者族 $\{P,B,O_s,O_t\}$。固定标签与截断 $H$ 空隙给相容扣减18。
若两个 $H$ 均命中核心、alpha 下截断均轻，两最大元截断给 $\xi\ge94+I$。
于是剩余预算至多40；再改用每侧只删一个最大元的同一个 $Q$，
其两 $H$ 都必须命中 $Q$。新空隙最多46个标签，使用较弱安全帽 $h_v\le R'+71$，
最终给 $44+2I+12\xi\le8\tau$，与 $\tau\le134-I$ 相差至少 $100+22I$。
一重 alpha、一个无核心命中的 $H$、两个都无命中的三种情形分别由
$100+2I\le0$、$\tau\ge91+I>62-I$、$144>134$ 排除。
正洞由五种可能低洞模式给 $\tau-\xi\ge65$，第一种情形已是
$\xi\ge86+I$ 与 $\tau\le138-I$ 的矛盾；其余三种仍有严格余量。

\subsection{五空的全部分支}
先用首项和两个见证星形证明两个 $O$ 均空，beta 与 $H$ 不能同侧活动；
随后证明整性、跨中心链和同侧链。每个配对的所有者必须来自实际贡献。
零洞预算 $\A\le76$。A0 保留 $P$ 的十三种布局由十一张实际仿射端点卡覆盖，
最小余量7；空 $P$ 的 $H/C_5$ 额外空单元给扣减13，
由 $2\tau-2\xi+7F\ge444$ 得 $483>482$。

alpha 额外空单元必须另证。初始扣减14并不意味着正常性。
先以通用核心命中端点和完整尾部目录得到扣减27；再区分两个物理侧容量，
在同一个实际加权截断上证明
\[
 |Q|+6F+3\tau-\xi\ge466,\qquad |Q|+6F+2\tau-\xi\ge458.
\]
第一物理侧给 $B\ge a$ 时上界559，而 $466+4\cdot27=574>559$，故 $B<a$、扣减31；
固定前缀计数给 $C<b$、扣减35；第二物理侧给 $D\ge a$ 时
$458+3\cdot35=563>559$，故 $D<a$、扣减38；再五个 LOW 标签给扣减42。
此时才可证明所有保留中心正常。长位移交叠和实际源边缘项使最终运输式给
\[
 884\le4I+4\tau-2\xi_s+884
 \le542+2(\Theta+B_s+B_t)\le882,
\]
矛盾。把实际边缘项提前放松为20会丢失此矛盾。

A1/B 的九个实际端点运输式分别有最小余量
\[
 811/2,\ 203/2,\ 258,\ 230,\ 91/2,\ 1105/2,\ 305,\ 415,\ 223.
\]
唯一未由这些运输式覆盖的空 $P$、额外空 $C_{5t}$ 族，使用先前完整的尾部六槽证明，
其相容扣减18给 $\tau+\delta\ge78>76-I$。该局部定理的整性、唯一 beta 贡献侧、
标签符号和尾部覆盖均在原报告中已证明，本次保持其精确假设。

正洞五空时 $\delta\le11$，先由 $K-M-\delta\ge58$ 排除 $O$ 与 beta/$H$ 活动。
扣减8后剩余至多3，所有原生普通端点正常。A0 空 $P$ 或 split beta 被扣减13或14排除；
A1 空 $P$ 的共同11个 HIGH 与18个 LOW 只触17个基线组，扣减12即矛盾。
其余 A0 十一卡最小余量307，A1/B 四个母式余量分别761、326、501、461。
至此零洞和正洞的全部五空情形闭合。

\section{覆盖、验证边界与复现}
按 $z=0$、$1\le z\le14$、$z\ge15$ 三段划分：最后一段整体排除；
零洞由 $Z_0,\ldots,Z_6$ 及 $j>6$ 覆盖；低正洞由 $P_0,\ldots,P_5$ 及 $j\ge6$ 覆盖。
结构中的低容量非本原、半步长例外在要求本原之前就已排除，因此无循环前提。
对三/四空和五空，验证器各枚举两分支全部 $2^{20}$ 个组件指派，
仅使用已在正文及附录证明的活动蕴含。五空合法指派分为
A0 32、同侧6、空P 10、对侧18、四端点2、split beta 2、B 32、B四端点8。
这证明有限布尔归类无遗漏；它不从有限采样推断实数集合的结构。

在仓库根目录运行：
\begin{verbatim}
python -I -S -B verification/upper83_complete/check.py --output .work/upper83.json
python -I -S -B -O verification/upper83_complete/check.py --output .work/upper83_O.json
\end{verbatim}
入口校验固定源摘要，重新运行原完整84基础、新结构、实际分类、低前缀、正洞、
全秩和五空尾部检查，再核对覆盖及负面控制；不接受输入中的历史 PASS 字段代替计算。
所需 Python 3.10+，无需第三方库、求解器或网络。
写出证明是定理的数学依据；程序验证其有限依赖、整数参数、精确恒等式、完整指派和数值证书。
程序通过不等于用形式化证明助手认证了所有无限论证。

以下附录为与验证器同目录保存的逐项证明说明，包含全部本次新增推导。
其中局部引理写明共同源和保留形式，是证明顺序中的接口；这些前提由正文的结构及分类提供。
旧报告中的“尚未证明83”等历史状态不被改写，本报告在此前局部结果上完成总接入。
\appendix
\clearpage
\hypertarget{size84-structural-reduction-and-all-high-hole-strata}{%
\section{Size84 structural reduction and all high-hole
strata}\label{size84-structural-reduction-and-all-high-hole-strata}}

Base cf26c50. These arguments, the common source interface and all local
terminals form the size84 exclusion in the companion report. General
dependencies: size85 foundation lemmas, source Theta \textless= 2d + 8,
all-rank kernel. Their proofs are cardinality independent; finite size85
rows are NOT imported.

\hypertarget{fresh-parameters-and-layer-identities}{%
\subsection{Fresh parameters and layer
identities}\label{fresh-parameters-and-layer-identities}}

n84, d76, e68, m67, a = 33-k, b = 34 + k, I = k(k + 1). Lambda = 1157 +
I; Omega = 612-I; C = 2m + 4e + 14-Lambda = -737-I. C + Omega = -125-2I.
G = 142. T \textgreater= max(48, z, 142-2z), T \textgreater= kappa +
142-3z. Unpaired bound2977, odd q \textless= 53, at least7 opposite
pairs. Provisional deficient sign: 1 + T68 = 2347\textgreater24*84 + 77
= 2093. Disappearing sign:1 + T67 = 2279\textgreater24*76 + 77 = 1901.
Uniform rhoP \textgreater= -8, rhoQ \textgreater= 0 gives k \textless=
24, a \textgreater= 9.

Once common core lattice exists, exact pre-G D1 bound is -125-2I + 2z (
+ 16 for a \textless= 14). Outer comparison Lambda\textgreater294 +
Omega-rho is strict even with rhoP = -8. For a, b \textgreater= 15,
maximum removal R \textless= 71 + rho + kappa + z. Restrict R(E)
\textless= 2R-1, giving fresh: F2-2F1 =
-251-2I-c-rho-mu-2kappa-2theta-2ell2-2g; F1-2F0 \textless= -127-2I +
rho-mu + 2z-c-2theta-ell1-g. After actual source layer overlap: C0-U0
\textgreater= 126 + 2I + c-2z; C1-V0 \textgreater= 250 + 2I + c + mu +
rho + 2z; c = 76j + delta. C0 = 1156 + I + tau.

\hypertarget{small-boundaries}{%
\subsection{Small boundaries}\label{small-boundaries}}

A larger AP Q yields z0 and common lattice by residue counting; strict
pre-layer comparisons then G142. At a9 Omega12, nonAP Q costs -8 +
53\textgreater12, and AP costs T \textgreater= 142\textgreater20. At a10
Omega60, nonAP Q has kappa \textless= 16\textless T10-1; strict inverse,
T48 + 52-8\textgreater60. AP costs142\textgreater68. At a11 Omega106,
nonAP Q has kappa \textless= 63\textless T11-1. Strict inverse, G142
yields z \textgreater= 40. Reuse universal effective witness (r56,
beta5, C132), hQ\textgreater10 = \textgreater muQ \textgreater= 189 =
\textgreater rhoQ \textgreater= 125\textgreater106 + 8. AP
costs142\textgreater114. At a12 Omega150: both nonAP, actual sufficient
inverse failure gives rhoQ \textless= 73. Endpoint repair beta3, q14,
j19, H6 (see check\_structure.py), obtains joint lattice. Pre-layer
strict. Then T \textgreater= 142-12 and rho \textgreater= -8 + 50 imply
cost172\textgreater150. In actual strict inverse branch kappa \textless=
108, D1 \textless= -817. G142 gives z \textgreater= 17; effective (55,
4, 119) = \textgreater rhoQ \textgreater= 113; T48 + 113-8 =
153\textgreater150. AP branches handled by general AP argument below
(their c caps are negative). At a13 Omega192: failed inverse R94 beta4,
q14, j23, H8; cost c \textless= 192-48-(142-16) = 18\textless76, w
\textless= 32, j0 obstruction. Strict inverse: if z \textgreater= 13,
effective (54, 6, 168) givesrhoQ153 and48 + 153-1 = 200\textgreater192.
Otherwise c \textless= 26\textless76, w48, j0. At a14 Omega232: failed
inverse R118 beta5, q14, j24, H10 gives c \textless= 62\textless76, w40.
Strict branch coarse pre-layer strict since kappa \textless= 184. If z
\textgreater= 15, independently certified primitive r53 h15 profile
gives rhoQ \textgreater= 269\textgreater232 (no D/G circularity).
Otherwise c \textless= 70\textless76, w56. j0 in both. So surviving
sizes15..52.

\hypertarget{ap-branches-and-low-c-lemma}{%
\subsection{AP branches and low-c
lemma}\label{ap-branches-and-low-c-lemma}}

Both-AP branches must be divided at a = 33. For12 \textless= a
\textless= 32, kappa \textless= Omega-rhoP\textless T\_b, so residue
counting on the larger AP gives one residue. The AP-plus-set formula
then has no clipped gap: z = 0, hQ = 0, hP = kappa. This holds for any
integer intrinsic step of P; no two-residue bound is needed here. A
nonprimitive P is excluded by the low-c argument below. In the only
remaining balanced case a33, b34, Omega612\textless2T33 = 1122. On the
common rational scale let p, q be coprime intrinsic AP steps. The two
residue counts force p, q \textless= 2, leaving(1, 1), (1, 2), (2, 1).
Unequal steps have z0 and kappa equal to one AP support length minus1,
hence at leastT33. Strict pre-G layers and G142 give c \textless=
612-142-T33\textless152, w0, and the low-c lemma excludes them. Thus the
surviving both-AP case has common unit step. When ONLY the smaller
support is AP, the larger nonAP floor is rhoQ \textgreater= b-5. The
correct bound is kappa \textless= Omega-rhoP-(b-5)\textless2T\_a,
checked in every row12 \textless= a \textless= 33. This is the premise
for its at-most-two residue alternatives, including the parity-fiber
calculation below. When ONLY larger Q is AP, use rhoP \textgreater=
n(a), so kappa \textless= Omega-n(a)\textless T\_b in every row. This
licenses the claimed one-residue branch without the false balanced
Omega\textless T\_b assertion. For a \textgreater= 12, one clipped
residue or unique off-scale generator costs T\_a + T\_(b-2) + 2-2b +
rhoP\textgreater Omega (rhoP -8 for12, -2 for13, 0 afterwards). Larger
AP or one-residue unclipped smaller AP gives z0. Nonprimitive small side
cap c \textless= Omega-142-T\_a-rhoP \textless= 85\textless152. Since
w0, j0/j1 are impossible by already proved zero-hole lemmas; their use
only needs common integral source, not individual primitivity.
Two-residue smaller AP: normalize U = 2{[}0, u-1{]}, least odd W, t
\textgreater= 2 odd Q generators. Exact mass floor muQ + kappa
\textgreater= u + (b-t)(b-t + 1) + tW + t(t-1)-W-T\_b. Against Omega +
2b-2-rhoP, finite minimization forces W \textless= 17, hence z = W-1
\textless= 16 and W\textless u. Exact parity-fiber identity gives hP, hQ
\textless= kappa. The general pre-G D1 comparison (-125-2I + 2z plus
small16) is already strict; no special half-step layer constant is
needed. G142 gives c \textless= 93\textless152, w \textless= 64. If
c\textless152, either j0 (contradicts w \textgreater= 2d + 8 = 160) or
j1 with delta \textless= 75. For j1, w \textless= 64, K \textgreater=
2d-1 = 151 makes leading active. The original one-hole actual owner
grammar (using raw first cap12 + eta \textless= 76, double17 + eta
\textless= 81, and active C7 d \textgreater= M + 5 \textgreater= 80)
gives the standard six-edge terminal and beta/O owner side. The raw caps
remain sufficient at w64: single cap12 + eta could76 when eta \textless=
w, so still \textless80; double \textless= 81 gives t \textless= 40.5.
Thus same pair arguments hold. Source cap tau \textless= 380 + w
\textless= 444. All-rank phi\_s + phi\_t cannot reach42 since min(Tp +
Tq, p + q = 42) = 462\textgreater444, so F \textgreater= 27. Six-edge
forest gives2X + G \textgreater= 162-delta \textgreater= 87. Owner
identity and beta/O cap give2X + G \textless= (17 + w)/2-2 \textless=
38.5 (use eta \textless= w; no localization needed). Contradiction. This
excludes every half-step branch. IMPORTANT: use eta \textless= w, not
w/2 before localization.

\hypertarget{both-nonap-inverse-repairs}{%
\subsection{Both nonAP inverse
repairs}\label{both-nonap-inverse-repairs}}

Actual sufficient test holds automatically on larger side. Small test
floor succeeds k0, 1; failures k2..21 use endpoint choices specified and
verified in check\_structure.py. Same proof as previous endpoint repair,
but q can be any \textgreater= 20 satisfying Tq-b-1\textgreater R,
allowing prefix floor n(q). Fresh main repaired rows k2..18 have H
\textless= 18; diameter margins positive. Pre-G D1 strict. Repaired Q is
primitive; if P nonprimitive its cap c \textless= 65\textless76, w
\textless= 72\textless160, so j0 obstruction excludes. Thus P primitive.
For STRICT inverse main branches, z \textless= 62 + I makes D1 strict.
If z \textgreater= 63 + I, both cores individually primitive already
(strict inverse hole bound). We can use the UNIVERSAL dense h49
certificate BEFORE D-integrality; it does not depend on G142. If one
final support splits disjointly, it costs rhoX \textgreater= T\_(r-2) +
2-2r, while other hasrho \textgreater= profile49(s). Across all38
oriented sizes15..52, cost + profile49(s) + (63 + I)-Omega \textgreater=
77. Hence both final translates intersect. Maximum-removal hX/2
sharpening gives D1 \textless= -125-2I\textless0. This replaces all
large-hole endpoint repairs and is noncircular. Original source outer
layer then aligns and intersects.

\hypertarget{complete-high-hole-stratum}{%
\subsection{Complete high-hole
stratum}\label{complete-high-hole-stratum}}

With sizes15..52 and primitive cores, paired universal floors over a +
b67: h15: (mu, rho\_elementary) = (418, 288); h17: (461, 331); h21:(503,
373); h49:(539, 409). Use independently verified profile table (all
relevant r15..52 present). Two-max source gives J + U \textless= 624 +
w-I-tau-c.~Outer layer gives2L \textless= 1530 + w-2I-2c-mu-rho-2z. Then
M \textless= 836 + H + (rho-mu)/2 + z-I-c; core central width W0 = 1156
+ I + mu-H satisfies W0-M \textgreater= -904 + 4I + 3c + 3(mu +
rho)/2-3z. Bands15..16, 17..20, 21..48 give positive margins107, 224,
266. Their I + c caps214, 179, 191 and H caps246, 219, 289 give
individual core mass minusH minima139, 166, 96. Tailz \textgreater= 49:
kappa \textless= 203-I-c, H \textless= 406-2I-2c, I + c \textless= 154,
so each coremass\textgreater H (compute minimum). If j \textgreater= 1,
c \textgreater= 76, then z \textless= 127-I-(c-76), width margin
\textgreater= 518 + 3c-3z \textgreater= 365 + 3I + 6(c-76)\textgreater0.
If j0, use unlocalized universal raw bound w \textgreater= 2M + 10 and w
\textless= 4z to getM \textless= 2z-5 \textless= 401; W0 \textgreater=
mu + 2rho-68 + 3I + 2c \textgreater= 1289. Hence same strip propagates
in all cases, w \textless= z and even. For j \textgreater= 1 source A
\textless= 380 + w but tau \textgreater= 418 + 2z, contradiction
(stronger in later bands). For j0 and z \textless= 48, w \textless= 48
contradicts no-empty w \textgreater= 160; in tail tau \textgreater= 539
+ 2z and A \textless= 456 + w \textless= 456 + z contradict. All
high-hole families gone without old first-gap H1 transport.

\clearpage
\hypertarget{low-hole-localization-and-the-one-hole-terminal}{%
\section{Low-hole localization and the one-hole
terminal}\label{low-hole-localization-and-the-one-hole-terminal}}

The common integral primitive source and1 \textless= z \textless= 14 are
hypotheses. The common source hypotheses are supplied by structure.md.

\hypertarget{a-simpler-localization-valid-for-every-j-1}{%
\subsection{A simpler localization, valid for every j \textgreater=
1}\label{a-simpler-localization-valid-for-every-j-1}}

Source bounds hold uniformly: A \textless= 456-76j + w, e\_v \textless=
621-76j + w-A, w \textless= 4z \textless= 56. Positive cores have mu
\textgreater= 61. Set H = tau-mu. Since min\_r(T\_r + I) = 385 and
max\_r(r-I) = 34,

\begin{verbatim}
C_v-H>=76j-w-10+delta>=10+delta,
W0=C0-2H>=822+76j-w+2I+delta,
M_v=R_v+e_v<=655-76j+w-I-delta.
\end{verbatim}

At j1, w56, W0 \textgreater= 842 and every one of the eight peeled
source entries is \textless= 635. Thus each translated central interval
joins as each source entry is adjoined, without using a large coarse
layer gap. All j \textgreater= 1 satisfy stronger bounds. The
common-strip theorem gives w, z even, w \textless= z \textless= 14,
equal target hole counts and boundary localization. This replaces the
more complicated j \textgreater= 3 localization in positive.md and also
supplies P1/P2.

\hypertarget{p1-reuse-the-sharp-one-hole-terminal}{%
\subsection{P1: reuse the sharp one-hole
terminal}\label{p1-reuse-the-sharp-one-hole-terminal}}

With one ordinary whole hole, the active own centres determine all outer
coordinates integrally without requiring leading activity. Each sign has
an active C6 or C7, so d\_v \textgreater= M + 5 \textgreater= 80;
moreover its leading maximum is positive. Thus the leading
cardinality/loss inequality and raw-gap lower guard are licensed even if
its full leading resource has no target hit (its entire K loss is still
in delta). Its shift is already integral.

If P were empty, an actual beta contributor and the opposite C5 give
t\_opposite\textgreater d\_beta \textgreater= 80, contradicting the
active O cap \textless= 19. Thus P has a contributor v, forcing
t\_v\textgreater d\_o \textgreater= 80 and making O\_v the unique hole.
Its own beta cannot contribute, because H\_v is active and beta/H would
give t\_v \textless= 16. Hence branch A has sole beta on u = o, with
beta\_u, H\_u, O\_u active and all alpha, C5, U, V, C7 retained. The
identical six-edge rooted proof from the zero-hole one-hole lemma gives

\begin{verbatim}
2t_u>=I+max(99-5xi+tau,27+2xi)-w>=51+I-w.
\end{verbatim}

The double raw cap gives2t\_u \textless= 17 + w/2. Their difference is
at least34-3w/2 \textgreater= 13. This closes P1 without first-gap or
two-cap subcases. Source bounds guarantee at least14 core entries on
each side, so the actual six-rank trim used in the displayed max exists.

\hypertarget{p2}{%
\subsection{P2}\label{p2}}

The former long first-gap route has been replaced by the shorter
complete shared proof in z2\_p2.md. It uses a32-unit prefix defect and
the same four beta/O cards as Z2. No eleven-value seed or P2 first-gap
census is needed.

\clearpage
\hypertarget{a-shared-two-hole-proof-z2-and-p2-at-size84}{%
\section{A shared two-hole proof: Z2 and P2 at
size84}\label{a-shared-two-hole-proof-z2-and-p2-at-size84}}

Common source hypotheses: d76, m68, r\_s + r\_t67, complete targets (Z2)
or localized positive targets with even w \textless= 14 and eta\_v = w/2
(P2). A = I + tau + delta \textless= 304 + w; R = 67 + xi; K-M
\textgreater= R-302-w + A. The source max bounds are e\_v \textless= 469
+ w-A and M\_v \textless= 503 + w \textless= 517 after localization (the
larger559 initial bound is also valid). Every integral ordinary centre
in{[}1-M\_v, S\_v{]} is active; the P2 proof of this property is
supplied below. For Z2 it follows simply from the actual source points0,
M\_v and S\_v \textgreater= M\_v.

\hypertarget{restore-leading-activity-in-z2}{%
\subsection{Restore leading activity in
Z2}\label{restore-leading-activity-in-z2}}

If leading v is empty, delta \textgreater= K and the source inequality
K-delta \textgreater= I + tau-8 gives I + tau \textless= 8, 31
\textless= r\_v \textless= 36 and u = delta-K \textless= 8-I-tau. Both
leadings cannot be empty because K \textgreater= 151\textgreater u. The
opposite leading has an integral shift. If at least two of the opposite
four own centres are holes, all four own centres on v are retained and
determine integral v-coordinates. Otherwise three opposite own centres
together with its integral leading shift determine all opposite
coordinates. Then an active C6\_v or C7\_v gives integral lambda\_v; if
both are holes all four own centres on v are retained. Thus lambda\_v is
integral in all cases. The usual two-witness residue argument, now with
both leading shifts integral, gives all outer coordinates integral; it
does not need a hit in both leadings.

Every other retained singleton has native loss \textless= u \textless=
8, so is normal: a BAD centre has at most max(r\_s, r\_t) + 4 \textless=
40 hits and loss \textgreater= 36. Since A \textgreater= K \textgreater=
151, e\_v \textless= 318. Thus S\_v-2M\_v \textgreater= T31 + 2*31 +
6-318 = 246\textgreater0. A leading centre in {[}1-2M\_v, S\_v{]} would
hit via0 or2M\_v. If lambda\_v\textgreater S\_v, at least two of its
four own centres remain retained and HIGH, each losing r\_v +
5\textgreater u points. Hence lambda\_v \textless= -2M\_v.

All four cross centres on v are then LOW. In A this removes three cross
singletons, impossible. In B the two exact holes must be C6\_v, C7\_v;
both paired resources have only the opposite contributor. All own
centres are retained. Normal V\_v has V\_v \textgreater= 72-r\_v-u. But
V\_v = lambda\_v + v4 \textless= X\_v-M\_v, so X\_v \textgreater= M\_v +
72-r\_v-u \textgreater= 76-u \textgreater= 68. Active C7\_o gives t\_o
\textgreater= 2X\_o + X\_v + g\_o + 1\textgreater69, contradicting
active O\_o's cap12. Thus both Z2 leadings are active. The same residue
argument proves outer integrality without any exceptional leading cases
remaining.

\hypertarget{p2-activity-and-a-minimal-positive-core-defect}{%
\subsection{P2 activity and a minimal positive-core
defect}\label{p2-activity-and-a-minimal-positive-core-defect}}

The common low-hole localization gives A \textless= 318 and18 \textless=
r\_v \textless= 49. K-delta \textgreater= -8-w + I + tau\textgreater0
using tau \textgreater= 61 + 2z and w \textless= z. Thus both leading
shifts are integral and the two-witness argument gives outer
integrality. Source side floors R\_v \textgreater= r\_v + 1 and tau\_v
\textgreater= r\_v-3 imply

\begin{verbatim}
S_v-H-R_o>=T_r+7r-606-2w+2I+2delta>=10+2delta>0.
\end{verbatim}

The quadratic minimum before subtracting606 + 2w is644. A target-hole
centre has r\_o \textgreater= 18 distinct opposite-core images in the
positive hull, more than eta\_v \textless= 7 missing targets. Nonholes
use zero; nonpositive centres use M\_v below the boundary strip, since
S\_v-M\_v \textgreater= C\_v\textgreater H. This proves the activity
interval used above.

Let f be the first missing absolute-core subset sum. The greedy rule
gives f\textless max core \textless= 352, since r\_v-I \textless= 34 and
I + tau \textless= 318. The side with at least34 core entries has a
prefix obtained by deleting its maximum, of mass at least T33 = 561.
That prefix still misses f, hence its subset sums are not a complete
interval. For n distinct positive integers with an incomplete subset-sum
interval, their mass excess over T\_n is at least n-1. Indeed, if the
first greedy failure follows k entries, all n-k subsequent rank
deviations are at least T\_(k-1) + 1 for k \textgreater= 1; their sum is
at least n-1. For k0 the bound is n.~The elementary inequality
(n-k)(T\_(k-1) + 1) \textgreater= n-1 holds for1 \textless=
k\textless n, first at n = k + 1 and then increases in n.~Applied to
this prefix, it proves

\begin{verbatim}
tau-xi>=32                       (positive case only).
\end{verbatim}

\hypertarget{no-two-caps}{%
\subsection{No two caps}\label{no-two-caps}}

The common four-top and star calculation gives 4I + 4tau + delta
\textless= 82 + 6w if both cap predicates hold. For P2 this contradicts
tau \textgreater= 61 + 2z, w \textless= z directly. For Z2 it gives tau
+ delta \textless= 82. If any C7\_v is active, the valid four-edge arm
gives t\_v \textgreater= 249-tau-delta \textgreater= 167\textgreater12.
If both C7's are empty, they are the exact holes and their lower C6's
are retained. Both C7's are HIGH, so the leading losses total at least
r\_s + r\_t + 160 = 227\textgreater82. Therefore no two caps exist in
either case. At least one O is empty. Here a cap means active O or
active beta, H.

\hypertarget{four-exact-betao-exclusions}{%
\subsection{Four exact beta/O
exclusions}\label{four-exact-betao-exclusions}}

Suppose beta\_v, O\_v are active. The sharp double cap says 2t\_v
\textless= 17 + w/2; the opposite O is the first hole. Let W be the
other whole hole. All listed native resources below are original,
distinct retained resources; every leading is kept at its full K
capacity. Roots explicitly retain their native loss in the forest.

\begin{enumerate}
\def\labelenumi{\arabic{enumi}.}
\tightlist
\item
  W is none of V\_v, C7\_v, C7\_o. Choose a full physical-v two-edge arm
  and the first edge of the other, rooted at leading v. The unused C7\_v
  has upper image t\_v-2X\_v-X\_o-g\_v \textless= t\_v-4, hence native
  loss \textgreater= 80-t\_v. This quantity is \textless= 76 whenever
  C7\_v is active; otherwise the assumed case is already impossible. The
  rooted forest gives t\_v \textgreater= K-M + 3F-delta + 80-t\_v. Since
  tau + 3F \textgreater= 198, this implies 2t\_v \textgreater= 43-w + I
  + xi. It contradicts the double cap because 26 + I + xi \textless=
  3w/2 \textless= 21 would be necessary.
\item
  W = V\_v. It is HIGH, since active leading implies
  V\_v\textgreater1-M\_v. Active beta gives
  alpha\_v\textgreater U\_v\textgreater V\_v. Thus the three retained
  alpha\_v, H\_v, O\_v each lose at least r\_v + 5. Use shortcut edges
  leading\_v--U\_v and leading\_o--C7\_o, root leading v, and the three
  unused losses. The full-core overlap F0 = 68-xi gives t\_v
  \textgreater= -84-w + I + tau-xi + 3r\_v \textgreater= 14-w +
  (tau-xi). In Z2 this is \textgreater= 14\textgreater8.5. In P2 it is
  \textgreater= 46-w \textgreater= 32\textgreater12 by the just-proved
  defect32 lemma.
\item
  W = C7\_v. It is HIGH above retained C6\_v. Use the retained
  three-edge forest leading\_v--V\_v and leading\_o--C7\_o--C6\_o, root
  leading v. The three unused U\_v, H\_v, O\_v lose at least r\_v + 5
  each. Hence t\_v \textgreater= -220-w + I + tau + xi + 3F + 3r\_v
  \textgreater= 76-w\textgreater12, using tau + 3F \textgreater= 198 and
  I + 3r\_v \textgreater= 98.
\item
  W = C7\_o. It is HIGH above C6\_o; its leading loses r\_o + 80 and the
  three unused V\_o, U\_o, H\_o lose r\_o + 5 each. Use edges
  leading\_v--V\_v and leading\_v--U\_v with both leading roots (the
  opposite is isolated). The total gap is2X\_v + g\_v. The full-core row
  gives t\_v \textgreater= -4-w + I + tau-xi + 4r\_o \textgreater=
  126-w\textgreater12, since I + 4r\_o \textgreater= 130.
\end{enumerate}

Thus simultaneous beta and O activity is impossible. The one-empty arm
rule is used only where it has at most one empty vertex per separated
arm; no capacity is restored to such a vertex.

\hypertarget{exhaustive-remaining-holes-and-a-common-terminal}{%
\subsection{Exhaustive remaining holes and a common
terminal}\label{exhaustive-remaining-holes-and-a-common-terminal}}

If both O's are holes, all other whole units are retained; B would give
both beta/H caps, so the branch is A and beta has a sole owner. If only
O\_v is missing, beta\_o cannot contribute because O\_o is active. In B
its singleton is the second hole, leaving beta\_v, H\_v active and two
caps. In A, either the entire beta pair is the second hole or its sole
contributor is v and H\_v must be the second hole. This exhausts exactly
\{O\_v, Bpair\}, \{O\_s, O\_t\}, \{O\_v, H\_v\} in A. All alpha, C5, U,
V, C7 singletons are retained. Choose the cap side u opposite the single
missing O, or the sole beta owner if both O's are missing. Actual
alpha\_o, C5\_u activity gives t\_u \textgreater= 2X + G + 2, hence2X +
G \textless= 10 + w/2.

The two six-edge retained trees on one F-optimal actual core give 2X + G
\textgreater= 6F-delta. The universal full-rank kernel tau + 6F
\textgreater= 378 then forces A \textgreater= 368 + I-w/2. Comparing A
\textless= 304 + w yields

\begin{verbatim}
64+I<=3w/2<=21,
\end{verbatim}

a contradiction. At w0 the margin is64. This closes Z2 and P2 under
their common-source hypotheses, without a first-gap table.

\clearpage
\hypertarget{positive-hole-activity-and-terminal-proofs}{%
\section{Positive-hole activity and terminal
proofs}\label{positive-hole-activity-and-terminal-proofs}}

Common integral primitive source and 1 \textless= z \textless= 14 are
supplied by structure.md. The uniform localization in
localization\_p12.md gives even w \textless= z \textless= 14 and eta\_v
= w/2 for every j \textgreater= 1. The exact finite replays are in
check\_positive.py. P0 is excluded by the unlocalized no-empty bound;
positive j \textgreater= 6 is impossible from A \textless= w \textless=
56\textless tau.

\hypertarget{one-dense-seed-and-a-ternary-prefix-census}{%
\subsection{One dense seed and a ternary prefix
census}\label{one-dense-seed-and-a-ternary-prefix-census}}

After localization and j \textgreater= 3, A \textless= 242, C0
\textless= 1398, r\_min \textgreater= 21, and each core entry \textless=
276. Six values missing from both sides in{[}1, 17{]} cost tau
\textgreater= 6*67 + 42-2*(12 + 13 + 14 + 15 + 16 + 17) =
270\textgreater242. Thus at least twelve distinct values occur there.
The retained exact6188-seed lemma gives {[}11, T\_B-11{]} in FS(B) and
at most five lower holes. Adjoin the other labelled core occurrences in
sorted order. For a current n \textgreater= 34, its total is at
least2T17 = 306, so x \textless= 276 \textless= T-21. For n \textless=
33 a failed join x\textgreater T-21 implies x \textgreater= 58 and at
least34 remaining entries, forcing C0 \textgreater= 78 + 34*58 =
2050\textgreater1398. Thus every join succeeds and all lower holes are
in{[}1, 10{]}.

The exact59049-pattern census enumerates multiplicities0, 1, 2 at 1,
\ldots, 10, then fills the remaining67-count slots from11 upwards, at
multiplicity at most2. It rigorously lower-bounds I + tau without
requiring an assignment to physical sides. Four lower holes cost at
least256; three cost at least195; the first gap2 costs at least99.
Consequently:

\begin{itemize}
\tightlist
\item
  j3: z \textless= 6, w \textless= 6, A \textless= 234.
\item
  j4: z \textless= 4, w \textless= 4, A \textless= 156; possible
  lower-hole sets are precisely a subset of \{ \{1\}, \{1, 2\}, \{1,
  3\}, \{2\}, \{3\} \}.
\item
  j5: the only lower hole is1. Hence z = w = 2, all positive core
  entries omit1, tau = 67 + sigma, and I + sigma + delta \textless= 11.
\end{itemize}

For j4, tau-xi \textgreater= 65. At first gap1 every rank deviation is
\textgreater= 1, so removing the two maxima leaves at least65 mass. At
first gap2 the low multiset is \{1\}; the other entries exceed2, giving
an even larger bound. At first gap3 it is \{1, 1\}; the remaining
entries exceed3, again giving at least65. This proves the bound directly
for every listed profile.

\hypertarget{activity-interval-j-3}{%
\subsection{Activity interval (j \textgreater=
3)}\label{activity-interval-j-3}}

The bounds above give S\_v-H-R\_o \textgreater= T\_r + 5r-480 + 2I +
2delta \textgreater= 113 + 2delta. Also S\_v-H-M\_v\textgreater0 from
the source side mass. A target-hole centre in{[}S\_v-H, S\_v{]} has at
least r\_o \textgreater= 21 distinct opposite-core images strictly
inside the target hull, more than eta\_v \textless= 7 holes. Thus it is
active. Other centres in{[}1, S\_v{]} use zero, and centres in
{[}1-M\_v, 0{]} use M\_v in the complete lower target strip. Therefore
every integral centre in{[}1-M\_v, S\_v{]} is active. An empty centre
above an active one is HIGH\textgreater S\_v. This supplies the
positive-case activity rules used by the zero-case classifiers without
assuming complete targets.

\hypertarget{fresh-two-cap-identity-and-j3j4-grammar}{%
\subsection{Fresh two-cap identity and j3/j4
grammar}\label{fresh-two-cap-identity-and-j3j4-grammar}}

The four-top tail has at most4*(76 + 6) + 5 = 333 values per sign: d\_s
+ d\_t + G + delta-2K \textgreater= 786-76j. The four full-core stars
give 2X + 4R + delta \textgreater= 540. With both caps t\_v \textless=
12 + eta\_v, put Y = 2M-2K + 24 + w + delta. Then G + 2X \textless= Y, G
\textgreater= 786-76j-Y, and Y \textless= 932-152j +
3w-2R-2I-2tau-delta. Hence, whenever the stars are licensed,

\begin{verbatim}
4I+4tau+delta+228j<=538+6w.
\end{verbatim}

This excludes j3 for w \textless= 6. For j4, before integrality or the
stars, the first two inequalities give2Y\textgreater482, but Y
\textless= 190 + 3w \textless= 202, so two cap predicates are already
impossible.

The cardinality-independent residue classifier in Appendix D now works
with fresh strict arm bounds. For j3 the two-witness argument
immediately gives integrality. For j4, two nonintegral lower coordinates
use four holes and leave both O active. With just one nonintegral lower
coordinate of residue theta and top residue b, the five residues H, O,
alpha, beta, C5\_opposite are b, b + theta, b-theta, b-2theta, -b. If b!
= 0 the actual branch-A/B hole counts again leave both O active. If b =
0 the exact holes are the two witnesses, O\_v, alpha\_v. The original
four integral-gap aggregate edges give at least248-tau-delta
\textgreater= 96-w \textgreater= 92; either active opposite O, C6 or
beta, O, C7 contradicts caps \textless= 14. A nonintegral top pair after
lower integrality costs four holes and the five-edge bound300-tau-delta
\textgreater= 148-w \textgreater= 144 contradicts the retained opposite
O, C6. Thus all outer coordinates are integral.

The separated-arm pointwise rule allows one empty vertex per arm with
zero capacity; failure to find an arm at j4 is exactly the four-hole
witness configuration already excluded. Therefore active C6 forces t
\textgreater= 301-tau-delta, and active C7 forces t \textgreater=
249-tau-delta. At j3, w \textless= 6 these are at least67 and15; at j4,
w \textless= 4 they are at least145 and93. They exclude O/C6, beta/H/C6
and beta/O/C7 simultaneous activity.

A HIGH cross makes leading loss at least r\_v + 80. At j3, empty C6
above active C5 leaves at least three higher native resources, each
losing r\_v + 5. The ten positive-positive fixed labels are HIGH and
compatible with all baseline groups. Thus A \textgreater= I + 4r + 105
\textgreater= 235, contradicting A \textless= 234. At j4 the C4 =
\textgreater C5 and C5 = \textgreater C6 implications have A
\textgreater= I + 3r + 90 \textgreater= 188\textgreater156. C6 =
\textgreater C7 uses one higher resource and the ten HIGH labels: A
\textgreater= I + 2r + 95 \textgreater= 161\textgreater156.

The literal owner identities and hole counts therefore give the same
complete j3 A0/A1 and j4 A-E/A-H/B tables as Appendix D. For j4's beta/O
exclusion, the three-edge rooted arm when V\_v is retained gives t
\textgreater= 119-w + I + xi; if V\_v is missing, its HIGH status and
three unused higher native losses give t \textgreater= 68-w + I + 3r\_v.
Both exceed the sharp double cap. Consequently both O are holes. The
remaining two-beta-hole B alternative has necessary row 2T\_r + 9r + 4I
+ 7tau/4 + 2tau\_v \textless= 1104 + 4w \textless= 1120, while its
integer quadratic minimum is1160. It follows from the actual HIGH-beta
master2S\_v + K + 3g\_v + 6X\_v \textless= M\_o + 1147 + w + delta and
the three-edge bound186-3tau/4-delta. The owner-side HIGH beta
contradicts 2d\_v \textless= M\_o-1 and actual P activity
d\_v\textgreater M. Thus neither extra beta-hole alternative is omitted.

\hypertarget{j3-terminal-bridge}{%
\subsection{j3 terminal bridge}\label{j3-terminal-bridge}}

The A0 geometric master in z34\_cards.md only changes1026 to1026 + w.
Use the relaxed common domain A \textless= 234, L \textgreater= max(8,
A-82), max(4, A-162) \textless= e\_v \textless= 399-A, and keep the
heavy thresholds70/72. The light-rank product276 still exceeds tau + eps
\textless= 226. The exact twelve minima at w6 are 162, 224, 285, 330,
126, 266, 422, 88, 158, 313, 41, 220, all positive.

For A1, the same actual side-union argument gives47 + w in place of47.
Its four Z upper bounds become212 + 3w, 288 + 3w, A + 52 + 2w, A + 128 +
2w. Their largest value is374, below Z \textgreater= 391.5. Thus all j3
positive-hole forms are excluded without individual first-gap terminal
cards.

\hypertarget{j4-terminal-bridge}{%
\subsection{j4 terminal bridge}\label{j4-terminal-bridge}}

Each of the seven Z4 actual masters in z34\_cards.md gains at most w on
its right. Use A \textless= 156, L \textgreater= max(8, A-4), max(4,
A-84) \textless= e\_v \textless= 321-A. Include every heavy flag allowed
by e + 72h \textless= 156 (even two in Card O if needed). The seven
finite minima are290, 65, 237, 151, 164, 283, 103, all positive.

For the exact ownerless family use z4\_ownerless.md with target hole
allowances. The initial four-arm bound is116 + I-w; hence X + g\_v
\textgreater= 57. Use the stronger active-C5 identity to obtain
d\_v-h\_v-X\_v \textgreater= X + G + h\_o-55-eta\_v\textgreater0 even at
w4. Thus the ten-label gap yields h\_v \textless= M + 11 + eta\_v. The
same e18 ledger works: V2\_s + V2\_t \textless= 20 + w-3X-G\textless0.
The core gap has at most30 + eta\_v positions, so h\_v \textless= R + 31
+ eta\_v with a core hit, or h\_v \textless= R\_v + M\_o + 30 + eta\_v
otherwise. Both H core-hit, light alphas imply xi \textgreater= 94 +
I-2w \textgreater= 86 + I. But the positive profile has tau-xi
\textgreater= 65 and tau \textless= 134 + w-I \textless= 138-I,
contradiction. One heavy alpha gives100 + 2I \textless= 3w \textless=
12. One H with no core hit gives tau \textgreater= 19 + u + I-2w
\textgreater= 83 + I, while delta \textgreater= 90 gives tau \textless=
66-I. Two non-core-hit H losses144 exceed eps \textless= 138. Thus the
ownerless positive family is excluded as well.

\hypertarget{j5-complete-terminal}{%
\subsection{j5 complete terminal}\label{j5-complete-terminal}}

The prefix census leaves just lower hole1. Thus z = w = 2, tau = 67 +
sigma and I + sigma + delta \textless= 11. All core entries omit1: xi\_s
= 1 + xi'\_s, xi\_t = 1 + xi'\_t, xi'\_s + xi'\_t \textless= sigma. In
particular r\_s, r\_t are between31 and36. Before outer integrality,
K-M-delta \textgreater= R-76 + I + tau \textgreater= 58. The leading
inequality rules out active O and active beta/H because their raw caps
are13 and10. The residue argument and literal owner identities of z5.md
now apply. A HIGH cross loses r\_v + 80\textgreater delta leading
points, supplying the same cross chains. Both O provide fixed credit8.
Every retained singleton or actual sole pair is normal: a BAD centre
loses at least72-max(r\_s, r\_t) \textgreater= 36, while eps \textless=
3. In particular all own chains and actual endpoint lists are licensed.

The source rows are freshly

\begin{verbatim}
S=2K+1074-delta, A=I+67+sigma+delta<=78,
L>=A+74, max(4,A-6)<=e_v<=243-A.
\end{verbatim}

Do not import the zero-hole strengthened bound L \textgreater= A + 78.
Use the actual full core q = 68, F = 68-xi\_s-xi\_t \textgreater= 63.
For A0 with empty P, the fixed ledgers for H/C5 extras have credit13 and
the alpha extra has credit14, each exceedingdelta11. Split beta has both
H HIGH and credit14, also impossible.

For A1 with empty P, the shared HIGH11 ledger from O\_s, O\_t, H\_s
combines with these18 LOW labels: C55, C56, C57, C58 on both signs; V5,
V6, V7, V8 on the inactive-beta sign t; F3, F4, F5 on both signs. They
touch17 baseline groups, so credit = 11 + 18-17 = 12\textgreater delta.
Pair1 is LOW, since a HIGH component would lose r\_v +
80\textgreater delta leading points. Beta\_t is LOW because one of H\_t,
alpha\_t is retained and normal. The F signs follow from
d\_v\textgreater g\_v. These signs hold for every possible final
companion hole. No inactive paired component is assigned native
capacity.

Retained-P A0 uses the same eleven exact affine cards as z5.md, with
normal endpoint errors zero. Their full-core minima, in A--K order,
are608, 452, 723, 428, 344, 504, 504, 307, 515, 597, 504. Card E uses
credit11 from HIGH H\_t; others usecredit8. Card D uses both actual
triples on this same full core, giving5q + 6F. For remaining A1/B, use
opposite/same/four-endpoint/branch-B physical masters from z5.md with
their right side increased by2. Even retaining their old nonnegative xi
error allowances and omitting the nonnegative side-mass terms, the
minima are 761, 326, 501, 461. check\_positive.py enumerates the entire
integer domain I + sigma + delta \textless= 11, xi\_v \textgreater= 1,
xi\_s + xi\_t \textless= 2 + sigma and the full peel polygon. This
closes every P5 form.

\clearpage
\hypertarget{fresh-geometric-derivations-for-z3z4-terminal-cards}{%
\section{Fresh geometric derivations for Z3/Z4 terminal
cards}\label{fresh-geometric-derivations-for-z3z4-terminal-cards}}

These are local terminals, assuming the actual retention forms and
common integral source. The classifiers in positive.md, z5.md and
check\_geometry.py provide the actual retention predicates separately
from these scalar terminals.

\hypertarget{z3-a0}{%
\subsection{Z3 A0}\label{z3-a0}}

S = 2K + 1224-delta, A \textless= 228, L \textgreater= A-76, max(4,
A-156) \textless= e\_v \textless= 393-A. The five weighted endpoints are
P\_s, 2alpha\_s, alpha\_t, C5\_t, H\_s. Their literal identity is H\_s =
P\_s + 2alpha\_s + alpha\_t + C5\_t + 4(M + X) + 2G. The endpoint
constant is5*72 + 71-3*67 = 230. The residual physical source mass is at
least T\_r + 4r + tau\_t + 4xi\_t + 10. With K \textless= M + R + 21 and
the six-edge retained forest, the exact master is

\begin{verbatim}
T_r+4r+tau_t+2L+4xi_t+12F<=1026-2e+delta+B_*.
\end{verbatim}

The common grouping rebates every endpoint's unused native loss except
the algebraically doubled alpha\_s. Consequently B\_* = 2b\_alpha\_s +
sum\_(other heavy Z)(b\_Z-k\_Z). Call other endpoints heavy at clip70;
alpha\_s's flag a is reserved for a peeled clip \textgreater= 72. If a =
0 its error \textless= xi\_s. All other light errors satisfy b
\textless= k: outside-hull clips have24 \textless= k \textless= 69, and
k(73-k) \textgreater= 276\textgreater220 \textgreater= tau + k. The
next-rank inequality (72-k)b \textless= tau proves the claim. Normal
endpoints have error0. Every heavy error is at most xi\_v + e\_v-4,
including core ranks70, 71. When a = 1, alpha\_s \textless= 0 and the
eleven LOW labels U5:8\_s, V5:8\_s, F3:5\_s touch nine groups. With the
eight O HIGHs, this gives compatible e10. Otherwise use e8.

Thus with n\_s in\{0, 1\}, n\_t in\{0, 1, 2, 3\},

\begin{verbatim}
B_*<=2xi_s+(2a+n_s)(e_s-4)+n_t(e_t-4)
      +n_s xi_s+n_t xi_t-70(n_s+n_t),
delta>=e+72a+70(n_s+n_t).
\end{verbatim}

These are precisely the twelve scalar systems exhaustively checked by
the full-rank verifier, with minimum margin65. The actual five resources
and the six-edge forest have been specified; no additional capacity is
assigned to an inactive component.

\hypertarget{z3-a1}{%
\subsection{Z3 A1}\label{z3-a1}}

The actual holes are O\_s, O\_t, H\_s; beta\_s and P\_t are sole
resources. The same twelve HIGH labels as in the former physical-union
proof give e12. Group leading s and its seven native resources: S\_s
\textgreater= K + 532-eps. Lower P\_t and upper H\_t are the only
endpoints. Let B = Gamma\_Pt + Gamma\_beta\_s and H = Gamma\_Ht
\textless= 75. The ten-edge forest and the two X\_s star edges give, on
one Q, Qgap + X\_s \textgreater= 6F-eps + B + H/2. The identity H\_t =
H\_s + P\_t + 2s4 + t4 + g\_s + g\_t + h\_t and HIGH H\_s imply S\_t
\textgreater= S\_s + M + M\_s + Qgap + X\_s + 77-Eclip-Eerr. As Eclip
\textless= B + H, the source sum1224 gives exactly

\begin{verbatim}
M+M_s+6F-H/2<=47+2delta+Eerr.
\end{verbatim}

The four endpoint flags p, h have Eerr \textless= xi\_s + xi\_t +
p(e\_t-4) + h(e\_s-4). Define Z = r\_s + 2I + 2tau + xi\_s + 6F-H/2. The
source peel bounds give, respectively, Z \textless= 212, 288, A + 52
\textless= 280, A + 128 \textless= 356. But2tau + 6F \textgreater= 396,
r\_s + 2I \textgreater= 33 and H \textless= 75 give Z \textgreater=
783/2 = 391.5. All endpoint positions are excluded with margin at
least71/2.

\hypertarget{z4-physical-capacities-and-masters}{%
\subsection{Z4 physical capacities and
masters}\label{z4-physical-capacities-and-masters}}

S = 2K + 1148-delta, A \textless= 152, L \textgreater= A, max(4, A-80)
\textless= e\_v \textless= 317-A. The generic rank constants are72-r\_v
and S\_v + r\_o-71. A native clip\textless72 has error \textless= xi\_v;
a heavier one has error \textless= xi\_v + e\_v-4 and costs at least72.
All heavy flags allowed by e + 72h \textless= 152 must be included, in
particular two flags when e8. No unproved assertion that at most one is
heavy is used.

The original endpoint/forest lists for F, S, EH, EA-own, EA-opp, EC5 are
retained literally. Physical t with seven native resources gives S\_s
\textless= K + 616-e; six gives S\_s \textless= K + 692-e. These are
union-deficiency bounds, not claims that their images are disjoint.
Direct substitution of the endpoint identities gives these fresh masters
(E is the actual weighted error, N the listed endpoints' native loss):

\begin{itemize}
\tightlist
\item
  F: T\_r + 4r + tau\_t + 2L + 4xi\_t + 12F \textless= 950 + delta-2e +
  E. The nominal constant is230. The six-edge forest avoids all five
  endpoints and Eclip \textless= 2N. There may be one empty ordinary
  vertex, charged by the separated-arm one-empty rule, with zero
  capacity.
\item
  S: 2M + 9F + 4e + 2N \textless= 539 + 3delta + E. Nominal constant153;
  physical t has six resources. The leading-t, six-edge and three-edge
  forest budgets rebate all four endpoints.
\item
  EH: R + 2L + xi\_t + e\_t + 10F + 3e + N-h \textless= 412 + 2delta +
  E, where h is the upper V\_s deficiency. Nominal constant225-r\_t;
  physical t has seven resources. The six-edge forest avoids all
  endpoints, while the four stars contain V\_s only. Drop N-h
  \textgreater= 0.
\item
  EA-own: L + 2e\_t + 2xi\_t + 10F + 3e + N \textless= 415 + 2delta + E.
  Use alpha\_s \textgreater= 1-a and physical t's seven resources. All
  four endpoints avoid both the eight-edge and two-star systems.
\item
  EA-opp: R + 2L + xi\_s + e\_s + 11F + 3e \textless= 487 + 2delta + E.
  Use alpha\_t \textgreater= 1-b; nominal constant159 + r\_t; physical t
  has six resources. The six-/five-edge forests avoid all weighted
  endpoints.
\item
  EC5: R + 2L + xi\_t + e\_t + 11F + 3e + N \textless= 488 + 2delta + E.
  Nominal constant225-r\_t; physical t has six resources. Its
  seven-/four-edge forests avoid the five endpoints.
\end{itemize}

Subtract baseline rank errors and each actual heavy increment e\_v-4.
Use N \textgreater= 72h where it occurs. The resulting systems are
exactly the six first Z4 cards in the full-rank verifier. Their margins
are 306, 81, 253, 163, 204, 295.

\hypertarget{z4-o-deriving-the-previously-conditional-constant994}{%
\subsection{Z4 O: deriving the previously conditional
constant994}\label{z4-o-deriving-the-previously-conditional-constant994}}

Its predicate is both O and H\_s empty, beta\_s native and P\_t sole,
all eight lower own/cross resources retained. The twelve HIGH labels
give e12. Put B = Gamma\_beta\_s + Gamma\_Pt, H = Gamma\_Vt. The actual
three forest rows are

\begin{verbatim}
X_s+g_s>=2F-(eps-B-H)/2,
X_s>=F-(eps-B-H)/2,
X_s+X_t>=2F-(eps-B)/2.
\end{verbatim}

Three, two and one copies give 6X\_s + X\_t + 3g\_s \textgreater=
10F-3eps + 3B + 5H/2. The endpoint identity V\_t = 2d\_s + beta\_s +
P\_t + 2a + b + g\_s, with twice leading s, has Eclip + 2ell\_s
\textless= 2eps-(B + H). Consequently

\begin{verbatim}
S_s+3g_s+4X_s+2a+b
<=1000+r_s+delta-2e-B-H+Eerr.
\end{verbatim}

Insert S\_s \textgreater= T\_r + tau\_s + 4R\_s + 6 + e\_s and the
forest row. This yields

\begin{verbatim}
T_r+4r+5xi_s+tau_s+R+L+2e_s+10F+60+4B+7H/2
<=994+4delta+Eerr.
\end{verbatim}

Thus994 is derived from the actual physical resources, rather than
extrapolated from1006. Baseline Eerr \textless= 2xi\_s + xi\_t. A heavy
beta\_s or P\_t costs4*72 = 288, a heavy upper V\_t costs(7/2)*72 = 252;
subtract its source e\_v-4. This is exactly the seventh full-rank Z4
system, with verified margin123.

The exact ownerless family, to which none of these cards applies, has
its separate proof in z4\_ownerless.md. The preceding activity lemmas
and the independent component census supply the exhaustive size84
classification; terminal arithmetic alone would not do so.

\clearpage
\hypertarget{size84-z4-ownerless-terminal}{%
\section{Size84 Z4 ownerless
terminal}\label{size84-z4-ownerless-terminal}}

Hypotheses: common integral source, complete targets, exact holes \{P,
Bpair, O\_s, O\_t\}, all fourteen singletons H, U, V, C5, C6, C7, alpha
retained, and integral outer coordinates. A = I + tau + delta \textless=
152, R = 67 + xi, L \textgreater= A, e\_v \textgreater= max(4, A-80),
e\_v \textless= 317-A. Neither empty pair is assigned a contributor.

\hypertarget{actual-fixed-credit-and-first-h-gaps}{%
\subsection{Actual fixed credit and first H
gaps}\label{actual-fixed-credit-and-first-h-gaps}}

The eight O HIGH labels give eps8 = delta-8. The four-edge arm on an
actual F-optimal core, with tau + 4F \textgreater= 260, gives 2X\_v +
X\_o + g\_v \textgreater= 4F-eps8 \textgreater= 116 + I. Thus active
C7\_v gives d\_v \textgreater= M + 117 + I and X + g\_v \textgreater=
59. Active alpha gives h\_v \textgreater= X\_v + 1. The original raw55
gap first gives h\_v \textless= M + 56= X + g\_v-55 \textgreater= 4.
Consequently exactly the same ten fixed candidates can enter (U\_v +
M\_v, H\_v-M\_o): C33, C34, C44, C15, C26, C27, C28, U2, U3, V2. Their
complete sign audit in Appendix D is valid with these strict margins. It
gives h\_v \textless= M + 11. Both empty paired centres are LOW. The
thirty-two LOW labels listed in that appendix touch25 groups; one
witness from each reflected pair V2, C35, C45 adds three LOWs in already
touched groups. For V2 use20-3X-G\textless0 since3X + G \textgreater=
232 + 2I. The untouched C25 group contributes another LOW witness. With
eight HIGHs this gives44 misses against26 baseline groups: e = 18.

The original core-gap list has37 candidates. Its initial bound h\_v
\textless= R\_v + M\_o + 37 removes all alpha peel contributions because
d\_v-h\_v-e\_v \textgreater= 80 + I\textgreater0. Also a-U3
\textgreater= X + g\_o + h\_o-36\textgreater0, where a = U\_v + R\_v. It
removes U3, U4, V2. The remaining30 candidates yield

\begin{verbatim}
h_v<=R+31 if H_v has a core hit,
h_v<=R_v+M_o+30 otherwise.
\end{verbatim}

In the second case the72 source points at or above -R\_o miss above the
complete target, so H\_v's native upper loss is at least72. The
whole-side source bound also gives S\_v\textgreater M: directly S\_v-M
\textgreater= T\_r + 4r-378 + I + delta \textgreater= 93 +
delta\textgreater0. This licenses the interval/core-hit statements
without hull gaps.

Let N be the sum of the two alpha deficiencies, k their lower clips, ell
the sum of the leading losses, and B the sum of the alpha rank errors.
Then alpha\_s + alpha\_t \textgreater= 77-k-B, N \textgreater= k. Write
eps = delta-18. For any actual common core with overlap f, the rooted
stars give

\begin{verbatim}
2X_v>=2f+N+ell-eps,   2X>=4f+N+ell-eps.
\end{verbatim}

An unused H loss u can additionally be rebated from the second row. For
light alpha clips (\textless72), B \textless= xi. At most one clip is
heavy, since eps \textless= 134\textless144.

\hypertarget{both-h-have-core-hits-both-alpha-clips-light}{%
\subsection{Both H have core hits, both alpha clips
light}\label{both-h-have-core-hits-both-alpha-clips-light}}

The H caps and alpha endpoints give L + X \textless= R-15 + k + B.
Combine the stars, L \textgreater= A, N \textgreater= k and k \textless=
eps to obtain

\begin{verbatim}
4f+2I+2tau+36+ell<=104+4xi.
\end{verbatim}

Delete two maxima on each core side: q = 64, f2 = 64-(tau-xi)/2. Its
substitution forces xi \textgreater= 94 + I. Thus tau \textgreater= 94 +
I, delta \textless= 58-2I, eps \textless= 40-2I.

Now delete just ONE maximum on each core side. Let the remaining maxima
be P'\_v = r\_v-1 + zeta\_v. Put zeta = zeta\_s + zeta\_t \textless=
tau-xi. This one actual Q has66 points, span R' = 65 + zeta, and overlap
f' = 66-zeta. Every H has a Q hit: otherwise at most the ten removed
source entries can hit, forcing native loss at least66\textgreater eps.
Every alpha clip is \textless= 40\textless71. Its next rank is strictly
below the core maximum, so its rank error is bounded by zeta\_v; hence B
\textless= zeta.

For completeness the Q-gap count is bounded independently. In (U\_v +
P'\_v, H\_v-P'\_o), only five positive non-Q entries from each of U, V,
alpha, five negative non-Q entries from H, and at most15 leading pair
labels can enter. Cross translates lie wholly below U\_v. A leading pair
above U\_v + P'\_v must use two of the five positive entries above
P'\_v. The same eleven fixed labels from the original truncated gap list
suffice: its exclusions use only P'\_v\textgreater0,
d\_v\textgreater h\_v, d\_v-o2 + P'\_o\textgreater0 and the already
established outer order. Thus there are at most15 + 5 + 15 + 11 = 46
labels. In particular the deliberately weaker safe cap h\_v \textless=
R' + 71 holds.

The two new H caps and the improved alpha errors imply

\begin{verbatim}
L+X<=2R'-R+65+k+zeta.
\end{verbatim}

Use the same rooted stars on Q and eliminate N, k as before. This
gives2I + 2tau + 2xi + 264 \textless= 220 + 10zeta, hence

\begin{verbatim}
44+2I+12xi<=8tau.
\end{verbatim}

But xi \textgreater= 94 + I and tau \textless= 134-I make the left side
minus right side at least100 + 22I\textgreater0. This excludes the
light/light case.

\hypertarget{other-halpha-statuses}{%
\subsection{Other H/alpha statuses}\label{other-halpha-statuses}}

If both H have core hits and alpha\_v has clip k\_v \textgreater= 72,
the opposite alpha endpoint and H cap give e\_o + X\_o \textless= 26 +
xi\_v + k\_o + zeta\_o, with zeta\_o \textless= xi\_o. Use the
individual rooted star on the one-trim support with f1 = 66-tau + xi, N
\textgreater= k\_v + k\_o and k\_o \textless= eps-k\_v. Then

\begin{verbatim}
2e_o+260+ell<=2(tau+delta)<=2A-2I.
\end{verbatim}

Since e\_o \textgreater= A-80, this would give100 + 2I + ell \textless=
0. The heavy alpha's unbounded peeled error was never estimated.

If exactly one H has no core hit, its unused loss u \textgreater= 72
leaves N \textless= eps-u \textless= 62, so both alpha clips are light.
Its truncated cap and the other H cap give e\_v + X \textless= R-16 + k
+ B. Rebate u in the rooted stars, use k \textless= eps-u and f1 =
66-tau + xi, obtaining

\begin{verbatim}
2e_v+198+2u+ell<=4tau+2delta.
\end{verbatim}

With e\_v \textgreater= A-80 this forces tau \textgreater= 19 + u + I
\textgreater= 91 + I. But delta \textgreater= 18 + u \textgreater= 90
gives tau \textless= 62-I, contradiction. If neither H has a core hit,
their two losses144 exceed eps \textless= 134. All cases are excluded
under the stated exact ownerless hypotheses.

\clearpage
\hypertarget{complete-local-z5-reduction-at-size84-common-source-assumed}{%
\section{Complete local Z5 reduction at size84 (common source
assumed)}\label{complete-local-z5-reduction-at-size84-common-source-assumed}}

The common integral source, complete targets and j = 5 are hypotheses.
The goal of this file is to close all actual Z5 geometries, not merely
scalar rows. Source notation is that of z5\_alpha.md. Use

\begin{verbatim}
S=2K+1072-delta, A=I+tau+delta<=76,
K-delta>=220+I+tau, K-M>=R-74+A,
max(4,A-4)<=e_v<=241-A, L>=A+76,
delta>=10 implies L>=A+78.
\end{verbatim}

Both leadings are active. The proof of S\_v\textgreater M in the
preceding tail report applies to every j5 source: r\_v \textgreater= 25,
R\_o \textless= 118, e\_o \textless= 241 give S\_v-M \textgreater= 325 +
75-118 + 6-241 = 47. Hence empty integral ordinary centres are either
LOW \textless= -M\_v or HIGH \textgreater= S\_v + M\_o + 1. A retained
translate with no core hit has at most four hits. After the eight O
credits each retained singleton, and each licensed sole pair, has at
least8 hits. Thus all of them have core hits and generic native endpoint
errors \textless= xi\_v (their clips are \textless= 68\textless72).

\hypertarget{the-base-holes-before-outer-integrality}{%
\subsection{The base holes before outer
integrality}\label{the-base-holes-before-outer-integrality}}

The sharp raw caps are t\_v \textless= 12 for active O\_v and t\_v
\textless= 9 for active beta\_v, H\_v. Here t\_v = d\_v-M. Leading
activity gives

\begin{verbatim}
t_v>=-7+I+tau+xi+2X_v+g_v.
\end{verbatim}

An active O therefore forces tau \textless= 19 and 2X\_v + g\_v
\textless= 19-I-tau-xi. If U\_v or C6\_o were active, q = X\_v + g\_v
would be integral. The two opposite leading/core shortcuts, including
the one-empty form, give

\begin{verbatim}
2q>=2(68-xi)-delta>=60+I+tau-2xi>=60+I-tau>=41.
\end{verbatim}

This contradicts q \textless= 19. Similarly active V\_v or C7\_o makes q
= X\_v integral and gives the same contradiction. Thus active O\_v
forces the four distinct singleton holes U\_v, V\_v, C6\_o, C7\_o. Two
active O's would force eight holes; so O\_o is empty. These are all five
holes, leaving V\_o, C7\_v retained. The corresponding two stars give
X\_o \textgreater= 21, whereas active C7\_v gives t\_v \textgreater=
2X\_v + X\_o + g\_v + 1\textgreater22. This contradicts t\_v \textless=
12. Both O's are empty.

If beta\_v, H\_v are active, the same argument with tau \textless= 16
forces four lower witness holes in addition to the two O holes. Thus
they cannot both be active. No lower coordinate integrality was assumed
in this argument: an active witness itself licenses the needed integer
displacement.

\hypertarget{integral-coordinates-and-the-literal-grammar}{%
\subsection{Integral coordinates and the literal
grammar}\label{integral-coordinates-and-the-literal-grammar}}

The unchanged four-coordinate identities in Appendix E of the retained
report now prove outer integrality. The needed numerical rules are
checked here rather than imported from size85. An empty cross above an
active cross is HIGH and makes the leading lose at least r\_v + 80,
which exceeds delta \textless= 76. Thus P\_v active = \textgreater{}
C5\_v = \textgreater{} C6\_v = \textgreater{} C7\_v. An empty integral
own centre above a retained own centre is HIGH; any retained larger
centre then loses72 points, more than epsilon \textless= 68. Thus
missing V forces U, H missing; missing U forces H missing.

Pair1 has a sole actual owner whenever retained, because P\_s + P\_t =
-(s2 + s3 + s4 + t2 + t3 + t4)\textless-3M whereas active centres
exceed-M\_v. The additional unchanged identities are

\begin{verbatim}
beta_s+beta_t+C5_s+C5_t=-2(s3+s4+t3+t4),
P_s+beta_s+C5_t=-2s3-2s4-t2-t3-t4.
\end{verbatim}

They exclude the indicated simultaneous activities. In A0 the beta pair
and both O's are holes. A nonintegral v3 forces U\_v, C6\_o holes, and a
nonintegral v4 forces V\_v, C7\_o holes. The remaining active H\_v,
alpha\_v and other lower own centre contradict respectively v3 =
H\_v-V\_v-alpha\_v or v4 = H\_v-U\_v-alpha\_v. Nonintegral v2 would cost
H\_v, alpha\_v, C5\_o. Thus all coordinates are integral. The chains
retain V\_s, V\_t, C7\_s, C7\_t. If P is retained the thirteen exact A0
templates are those listed in the full-rank checker (eleven endpoint
cards with the three orientations of A). If P is empty, the final hole,
up to reflection, is H\_s, alpha\_s or C5\_s.

For split beta, the identities force exactly O\_s, O\_t, H\_s, H\_t,
C5\_t as holes, with sole P owner s. The retained lower own centres and
alphas prove integrality. The split-beta transport below excludes this
geometry using only the whole beta capacity. For sole beta\_s, base
holes are O\_s, H\_s, O\_t. A nonintegral lower coordinate again costs
two extra holes. On t, H\_t, alpha\_t and the other lower own centre
contradict it. On s, alpha\_s, beta\_s make d\_s, s3 + s4, s2 integral
and contradict it. Nonintegral t2 costs three holes H\_t, alpha\_t,
C5\_s. The chains retain the four extreme leaves. If P is empty, its
final companion is precisely alpha\_s, H\_t, alpha\_t, C5\_s, C5\_t or
U\_s. If P is retained, it is either opposite beta's owner, the same
owner with alpha\_s retained, or the same owner with alpha\_s, C5\_t
empty.

In B, the strict beta-sum sign forces C5's pair retained with sole owner
s. Any retained P has the same owner. The beta/H exclusions cost one
hole on each sign in addition to the two O holes. The same two-witness
identities and own/cross chains prove integrality and extreme-leaf
retention. The remaining alternatives are alpha\_s retained, or alpha\_s
the only extra hole; in the latter case P, alpha\_t and all eight lower
own/cross resources are retained. No empty pair receives a nominal
owner.

\hypertarget{common-endpoint-and-forest-calculus}{%
\subsection{Common endpoint and forest
calculus}\label{common-endpoint-and-forest-calculus}}

Write eps = delta-e. Every sole native resource has lower and upper
bounds 72-r\_v-l-xi\_v and S\_v + r\_o-71 + u + xi\_o. Distinct original
resource deficiencies, or deficiencies of disjoint actual groups, sum to
\textless= eps. Native endpoint weights are recorded explicitly below.
For one actual Q of size q and overlap F, weighted stars and ordinary
gap pairs give

\begin{verbatim}
cX_s+dX_t >=(c+d)F-max(c,d)eps/2,
cg_s+dg_t >=(c+d)F-max(c,d)eps/2.
\end{verbatim}

These are separate budget consequences. Every group pair has at most one
empty member; the full intersection in that case is paid solely by the
retained partner's off-target points. A retained own triple gives g + h
\textgreater= q + F-Psi. It shares its actual deficiency with all of its
subgroups. No subtree or subgroup is allocated twice.

\hypertarget{a0}{%
\subsection{A0}\label{a0}}

Retained P uses the thirteen actual templates and eleven exact affine
cards A--K in \texttt{upper83\_fullrank/check.py}. That checker
reconstructs the endpoint identities at d76 and proves margins 321, 141,
324, 59, 7, 167, 167, 25, 203, 257, 167. Their source, activity,
core-hit and sole-ownership predicates have now been supplied above. For
E the empty H\_t is HIGH above retained U\_t, providing credit11; the
other cards use credit8. Card D uses both actual triples and one common
support optimizing 5q + 6F.

When both pairs are empty and the extra hole is H\_s or C5\_s, the
initial fixed ledger gives credit13 (HIGH11/LOW22/touched20 or
HIGH8/LOW25/touched20). It needs only LOW paired centres and the leading
inequalities; it does not require normal endpoints. The empty pairs are
LOW because retained H or alpha lies above each beta, and a HIGH cross
would lose r\_v + 80 leading points. Put B = alpha\_t, D = C5\_t. The
physical six-resource side gives S\_t \textless= K + 616-e. From H\_t =
2(a + b) + G + alpha\_s + B + D, one generic lower endpoint and the
actual B, D packing inequality give

\begin{verbatim}
M+2X+G+F <= R+484-e+u_Ht+l_alpha_s+2Gamma_BD+2xi.
\end{verbatim}

Rebate those three disjoint groups in the stars and ordinary pairs.
Their maximum coefficient is two, so M + 7F \textless= R + 484-e + 2eps
+ 2xi. Using L \textgreater= A + 78 and A \textless= 76 yields

\begin{verbatim}
2I+2tau-2xi+7F+3e<=482.
\end{verbatim}

The full-rank kernel 2tau-2xi + 7F \textgreater= 444 makes the left side
at least483. For the remaining alpha\_s hole, \texttt{z5\_alpha.md}
proves the full ordered
credit14-\textgreater27-\textgreater31-\textgreater35-\textgreater38-\textgreater42
path and the final884\textgreater882 contradiction.

\hypertarget{a1b-transports-reconstructed-at-d76}{%
\subsection{A1/B transports reconstructed at
d76}\label{a1b-transports-reconstructed-at-d76}}

The fresh physical inequalities are listed first. Every centre identity
below follows from a = M\_s + X\_s, b = M\_t + X\_t and the definitions;
the endpoint lists specify the actual resource accounting.

\begin{itemize}
\tightlist
\item
  Opposite owners beta\_s, P\_t: use V\_t = 2d\_s + beta\_s + P\_t + 2a
  + b + g\_s, upper V\_t, lower beta\_s, P\_t, and twice leading s. This
  gives S\_s + 3g\_s + 4X\_s + 2a + b \textless= 924 + r\_s + delta-2e +
  xi + xi\_s.
\item
  Same owners with alpha\_s retained: the beta-sum identity and
  alpha\_t\textgreater-b give alpha\_s-P\_s\textgreater beta\_s + d\_t +
  2a + b + G. The upper alpha\_s, lower P\_s, beta\_s and leading t give
  S\_t + M\_s + 2X\_s + 3X\_t + g\_s + 2g\_t \textless= R + 1012-r\_t-e
  + 2xi\_s + xi\_t.
\item
  Four endpoints: U\_s = P\_s + alpha\_t + C6\_t + 2a + 4b + 2G. The
  upper U\_s and three lower endpoints give S\_t + 2M\_t + 2X\_s + 4X\_t
  + 2G \textless= 2R + 894 + 2r\_t-e + xi\_s + 3xi\_t.
\item
  HIGH alpha\_s: use alpha\_s + C5\_t = d\_t-2(a + b)-G, d\_t =
  H\_t-alpha\_t, and alpha\_s \textgreater= S\_s + R\_t + 1. The three
  retained endpoints C5\_t, H\_t, alpha\_t give 2S\_s + 2X + G
  \textless= 2R + 965-e + xi. Here e14 includes both O's, HIGH H\_s and
  HIGH alpha\_s. The alpha HIGH guard follows from d\_s-R\_s
  \textgreater= 220 + I + tau-R\_s\textgreater0 and the core hit of
  beta\_s.
\item
  Empty P, extra C5\_s or U\_s: sum upper alpha\_s, H\_t minus lower
  2beta\_s, alpha\_t, C5\_t. The literal expression is 2d\_s + 2a + 2b +
  G. Twice leading s gives 2M + 6X\_s + 2X\_t + 3g\_s + g\_t \textless=
  843 + delta-2e + 3xi.
\item
  Empty P, extra H\_t: upper alpha\_s, V\_t minus lower 2alpha\_t,
  2C5\_s, 2beta\_s, C5\_t equals6a + 5b + 3g\_s + 2g\_t. Hence 4M\_s +
  3M\_t + 6X\_s + 5X\_t + 3g\_s + 2g\_t \textless= 2R + 736 + r\_s +
  delta-2e + 4xi + xi\_s.
\item
  Empty P, extra alpha\_t: upper alpha\_s, H\_t minus lower 2beta\_s,
  C5\_t, 2C5\_s equals6a + 3b + 3G + 3h\_t. Hence 4M\_s + M\_t + 6X\_s +
  3X\_t + 3G + 3h\_t \textless= 2R + 746 + 3r\_s + delta-2e + 5xi\_s +
  2xi\_t. Use the H\_t, U\_t, V\_t triple, two g\_s pairs and separate
  stars, obtaining3q + 15F-6eps for the gap terms.
\item
  Split beta: its entire hit count N has N \textgreater= 76-delta and
  beta\_s + beta\_t \textgreater= N-M. Use upper alpha\_s, lower P\_s
  and twice leading t in alpha\_s-P\_s = beta\_s + beta\_t + 2d\_t + 2a
  + 2b + G. Conservatively, S\_t + M + 2X\_s + 6X\_t + g\_s + 3g\_t
  \textless= 920 + 2delta + xi. Both H's are HIGH, so credit14 is
  compatible. N is charged only to the one whole pair, never to either
  component separately.
\item
  Branch B, alpha\_s retained: alpha\_s-C5\_s\textgreater d\_t + 2a + 2b
  + G. Two endpoints and leading t give S\_t + M + 2X\_s + 4X\_t + g\_s
  + 2g\_t\textless R + 1017-e + xi.
\end{itemize}

Replace each S\_v by T\_(r\_v) + tau\_v + 4R\_v + 6 + e\_v. The exact
checker uses the full source peel polygon and all-rank inverse mass
profiles; the minima of left minus right are respectively 811/2, 203/2,
258, 230, 91/2, 1105/2, 305, 415, 223. The two endpoint sides retain
physical names in these calculations; the scalar quadratic index ranges
over both orientations.

These nine exclusions route all non-tail A1/B alternatives. HIGH-alpha
handles extra alpha\_s; the next three rows handle C5\_s/U\_s, H\_t,
alpha\_t. The exact final empty-P C5\_t family is already fully proved
and pinned by \texttt{reports/upper83\_tail\_20261019.tex} and
\texttt{verification/upper83\_tail.py}. With P retained, use opposite
owners, same owners with alpha\_s retained, or the four-endpoint row if
alpha\_s is empty. Branch B uses its own row when alpha\_s is retained
and the four-endpoint row otherwise. Together with A0 this closes every
Z5 geometry under the stated common source hypotheses. The common-source
and positive-hole obligations are discharged separately in structure.md
and positive.md.

\clearpage
\hypertarget{size-84-z5-alpha-order-before-normality}{%
\section{Size-84 Z5 alpha: order before
normality}\label{size-84-z5-alpha-order-before-normality}}

Local theorem, assuming the common integral source and the exact
Branch-A whole holes \texttt{\{P,\ Bpair,\ O\_s,\ O\_t,\ alpha\_s\}}.
This file does not assert the global classification. All thirteen other
singleton resources are retained. Notation follows the full-rank report:
d = 76, m = 68, r\_s + r\_t = 67, I = (r\_s-33)(r\_s-34), A = I + tau +
delta \textless= 76, R = 67 + xi, M\_v = R\_v + e\_v, a = M\_s + X\_s, b
= M\_t + X\_t, G = g\_s + g\_t. Write B = alpha\_t, C = C5\_s, D =
C5\_t. In particular B is a centre here, not the empty beta pair or a
source fringe count.

\hypertarget{the-initial-tail-credit-does-not-need-normality}{%
\subsection{The initial tail credit does not need
normality}\label{the-initial-tail-credit-does-not-need-normality}}

The generic Z5 core-hit argument gives each retained centre Z at least
one core hit. The native rank endpoints are Z \textgreater=
72-r\_v-l\_Z-xi\_v and Z \textless= S\_v + r\_o-71 + u\_Z + xi\_o
whenever the selected native clip is below72. Also Z \textgreater=
1-R\_v is always available from a core hit. The empty alpha\_s and both
empty beta components are LOW. The original initial label inventory
therefore gives credit14.

With epsilon = delta-14, the full-rank stars and ordinary pairs give
X\_v, g\_v \textgreater= F-epsilon/2. Since tau + 2F \textgreater= 134
(min(T\_p-2p) = -1 on each side), X\_v, g\_v \textgreater= 36 + I/2,
hence at least36. The retained \{H\_s, U\_s\} pair gives h\_s
\textgreater= F-epsilon \textgreater= 6 + I, using tau + F \textgreater=
68. LOW alpha\_s gives h\_s \textless= X\_s. K-delta \textgreater= 75 +
I + tau gives d\_v \textgreater= 75 + I + tau + g\_v + 2X\_v.

These inequalities license the complete32-type tail inventory from
\texttt{upper84/audit/verification\_program/proof/report.tex} (the
paragraph ``The alpha\_s hole: the fixed-credit argument''). Here are
its size-specific strict checks. On s, six beta types lie below
d\_s-X\_s. The remaining 26 types give S\_s \textless= d\_s + 26. The
three labels C16\_s, C17\_s, C18\_s exceed that cap since C + s3 + h\_s
\textgreater= 1 + e\_s + X\_s + g\_s + h\_s \textgreater=
83\textgreater26. At d\_s + 1, d\_s + 2 every remaining type except the
sole score O\_s-t1 is either \textless= d\_s or \textgreater= d\_s + 6;
use h\_s \textgreater= 6, g\_s \textgreater= 36 and
d\_s\textgreater g\_s + 2h\_s. Thus S\_s \textless= d\_s + 1. Above H\_s
+ M\_s the only possible types are d\_s-g\_s, d\_s, O\_s-t1; they cannot
cover three consecutive integers. Therefore S\_s \textless= H\_s + M\_s
+ 2. On t, sixteen types lie below H\_t + M\_t and all nonfixed
resources end there, so S\_t \textless= H\_t + M\_t + 16. X\_t
\textgreater= 36 makes C22\_t, C23\_t, C24\_t HIGH; D + M\_t + 2X\_t +
g\_t \textgreater= 1 + e\_t + 2X\_t + g\_t \textgreater=
113\textgreater16 makes C16\_t, C17\_t, C18\_t HIGH. The s side gives
the analogous six labels and C33\_s. Thus HIGH21, LOW29, touched23 gives
credit27, before normality.

\hypertarget{separate-the-two-order-tests}{%
\subsection{Separate the two order
tests}\label{separate-the-two-order-tests}}

At credit e \textgreater= 27, epsilon = delta-e \textless= 49. All clips
are below72, so generic errors \textless= xi\_v apply even to centres
outside {[}1, S\_v{]}. The source lemma gives L \textgreater= A + 78
since delta \textgreater= 27. Grouping whole actual targets gives S\_s
\textless= K + 540-e and S\_t \textless= K + 616-e. Consequently the
exact identities H\_s = a + 2b + G + h\_s + B + C and H\_t = a + 2b + G
+ h\_s + B + D imply

\begin{verbatim}
M_t+X_s+2X_t+G+h_s+B+C <= R+490+r_t-e+u_Hs+xi_t,
M_t+X_s+2X_t+G+h_s+B+D <= R+566+r_s-e+u_Ht+xi_s.
\end{verbatim}

Choose ONE actual truncated core maximizing q + 6F. The weighted inverse
staircase gives, throughout tau \textless= 49,

\begin{verbatim}
q+6F+3tau-xi >=466,
q+6F+2tau-xi >=458.
\end{verbatim}

The checker reconstructs these minima using integer mass constraints.
The actual triple \{H\_s, U\_s, V\_s\} and the two g\_t pairs give G +
h\_s \textgreater= q + 2F-epsilon; stars give2X \textgreater=
4F-epsilon.

If B \textgreater= a, lower C generically in the FIRST physical row. As
r\_s + r\_t = 67,

\begin{verbatim}
M+2X+G+h_s <= R+485-e+epsilon+xi.
\end{verbatim}

Combining the preceding bounds, L \textgreater= A + 78 and delta
\textless= 76-I-tau gives

\begin{verbatim}
q+6F+3I+3tau-xi+4e <=559.
\end{verbatim}

At e = 27 its left side is at least466 + 108 = 574\textgreater559. Hence
B\textless a. This conclusion uses the smaller physical-s capacity;
replacing both physical rows by the larger one loses the required
margin.

The six new LOW labels U5\_t, U6\_t, U7\_t, U8\_t, V3\_t, V4\_t meet two
new groups. For V3\_t, V4\_t use D \textgreater= 1-R\_t and
b\textgreater R\_t; positivity of D is unnecessary. Credit is now31.

Independently C\textless b. If C \textgreater= b, it is the smallest
retained ordinary s centre and the leading starts at C + 2X\_t + B +
g\_t\textgreater C. For the strict sign, already at credit27, 2X\_t +
g\_t \textgreater= 127.5 whereas R\_t \textless= R \textless= 116, and B
\textgreater= 1-R\_t. Thus {[}1, C-M\_t-1{]} is fixed-only. The same
credit27 stars give X\_t \textgreater= 42.5, so it contains at least42
integers. But the s-side fixed inventory has11 HIGH and18 LOW among55
labels, leaving at most26. Contradiction. C\textless b gives C66\_s,
C67\_s, C68\_s, C35\_s LOW; for C35\_s = C-D-2b-g\_t-h\_s use D
\textgreater= 1-R\_t. No new group is touched. Credit is35.

Finally suppose D \textgreater= a. Lower B generically in the SECOND
physical row:

\begin{verbatim}
M+2X+G+h_s <= R+561-e+u_Ht+l_B+xi.
\end{verbatim}

Let N = Gamma\_Ht + Gamma\_B. The triple, g\_t pairs and stars used
above avoid both resources, so their respective budgets may each be
epsilon-N. As u\_Ht + l\_B \textless= N, this gives

\begin{verbatim}
L+q+6F <=561-e+2epsilon-N+xi <=561-e+2epsilon+xi,
q+6F+2I+2tau-xi+3e <=559.
\end{verbatim}

At e35 the latter left side is at least458 + 105 = 563\textgreater559.
Hence D\textless a. The three labels C66\_t, C67\_t, C68\_t add credit3,
giving38. C25\_s, C25\_t, C35\_t, C45\_s, C45\_t add credit4, giving42.
Their signs use C\textless b, D= 1-R\_s, D \textgreater= 1-R\_t, not
positivity. The two C25 labels touch the sole remaining baseline group.

\hypertarget{normality-and-terminal}{%
\subsection{Normality and terminal}\label{normality-and-terminal}}

Only now exclude BAD centres. Put r\_max = 34 + k, I = k(k + 1). A BAD
retained singleton loses at least72-r\_max = 38-k, whereas epsilon
\textless= 34-I-tau. Their difference is k\^{}2 + 4 + tau\textgreater0.
Thus all retained singleton centres are normal.

The exact transport of the full-rank report now applies with the actual
source fringe count B\_fr = B\_s + B\_t:

\begin{verbatim}
2M_t+2q+11F <=2R+778+2B_fr+2r_t-delta+5(delta-42).
\end{verbatim}

For completeness its eleven distinct singleton groups are \{B, D\},
\{H\_s, U\_s, V\_s\}, \{U\_t, V\_t\}, \{C6\_s, C7\_s\}, \{C\}, \{H\_t\};
the separately valid weighted stars use the same actual core. The native
losses are not counted as disjoint between these two separate
inequalities. At credit42 all required gaps first satisfy \textgreater=
tau. The boundary-safe complete-core intersection raises them to
\textgreater tau. The complete-core long-gap intersection and
three-translate union lemma therefore permits q = F = 68 simultaneously
in those groups. Use e\_t \textgreater= Theta-d-12 and Theta = 84 + A,
then retain the actual fringe in K \textless= M + R + 1 + B\_fr. The
transport becomes

\begin{verbatim}
4I+4tau-2xi_s+13*68 <=542+2(Theta+B_fr)<=882.
\end{verbatim}

The left side is at least884, contradiction. Thus this exact local
alpha-hole family is excluded, without an early normality assumption.

\end{document}
